% Môn: Vật lý
% Lớp: 10
% Chương: Chương III. Động lực học
% Bài: L10C3 Bài 18. Lực ma sát
% Số câu: 162
% App đọc file này trực tiếp — sửa rồi Commit trên GitHub, bấm Đồng bộ GitHub .tex

% L10C3 Bài 18. Lực ma sát

% ===== Câu 1 =====
% ID: AIQ_L10C3_FULL_0811
% Mức: NB
\dangbt{Nhận biết các loại lực ma sát và điều kiện xuất hiện}
\begin{ex}
% ID: AIQ_L10C3_FULL_0811
% Mức: NB
Một vật trượt trên mặt phẳng ngang. Áp lực lên mặt là 20 $N$, hệ số ma sát trượt 0,2. Tính lực ma sát trượt.
\choice
{\True 4 $N$}
{5 $N$}
{3 $N$}
{8 $N$}
\loigiai{
Chọn A. F_mst=μN=0.2·20=4 $N$, ngược chiều chuyển động.
}
\end{ex}

% ===== Câu 2 =====
% ID: AIQ_L10C3_FULL_0812
% Mức: NB
\begin{ex}
% ID: AIQ_L10C3_FULL_0812
% Mức: NB
Một vật trượt trên mặt phẳng ngang. Áp lực lên mặt là 30 $N$, hệ số ma sát trượt 0,3. Tính lực ma sát trượt.
\choice
{9 $N$}
{\True 9 $N$}
{8 $N$}
{18 $N$}
\loigiai{
Chọn B. F_mst=μN=0.3·30=9 $N$, ngược chiều chuyển động.
}
\end{ex}

% ===== Câu 3 =====
% ID: AIQ_L10C3_FULL_0813
% Mức: NB
\begin{ex}
% ID: AIQ_L10C3_FULL_0813
% Mức: NB
Một vật trượt trên mặt phẳng ngang. Áp lực lên mặt là 40 $N$, hệ số ma sát trượt 0,4. Tính lực ma sát trượt.
\choice
{16 $N$}
{17 $N$}
{\True 16 $N$}
{32 $N$}
\loigiai{
Chọn C. F_mst=μN=0.4·40=16 $N$, ngược chiều chuyển động.
}
\end{ex}

% ===== Câu 4 =====
% ID: AIQ_L10C3_FULL_0814
% Mức: NB
\begin{ex}
% ID: AIQ_L10C3_FULL_0814
% Mức: NB
Một vật trượt trên mặt phẳng ngang. Áp lực lên mặt là 50 $N$, hệ số ma sát trượt 0,5. Tính lực ma sát trượt.
\choice
{25 $N$}
{26 $N$}
{24 $N$}
{\True 25 $N$}
\loigiai{
Chọn D. F_mst=μN=0.5·50=25 $N$, ngược chiều chuyển động.
}
\end{ex}

% ===== Câu 5 =====
% ID: AIQ_L10C3_FULL_0815
% Mức: TH
\begin{ex}
% ID: AIQ_L10C3_FULL_0815
% Mức: TH
Một vật trượt trên mặt phẳng ngang. Áp lực lên mặt là 60 $N$, hệ số ma sát trượt 0,2. Tính lực ma sát trượt.
\choice
{\True 12 $N$}
{13 $N$}
{11 $N$}
{24 $N$}
\loigiai{
Chọn A. F_mst=μN=0.2·60=12 $N$, ngược chiều chuyển động.
}
\end{ex}

% ===== Câu 6 =====
% ID: AIQ_L10C3_FULL_0816
% Mức: TH
\begin{ex}
% ID: AIQ_L10C3_FULL_0816
% Mức: TH
Một vật trượt trên mặt phẳng ngang. Áp lực lên mặt là 20 $N$, hệ số ma sát trượt 0,3. Tính lực ma sát trượt.
\choice
{6 $N$}
{\True 6 $N$}
{5 $N$}
{12 $N$}
\loigiai{
Chọn B. F_mst=μN=0.3·20=6 $N$, ngược chiều chuyển động.
}
\end{ex}

% ===== Câu 7 =====
% ID: AIQ_L10C3_FULL_0817
% Mức: TH
\begin{ex}
% ID: AIQ_L10C3_FULL_0817
% Mức: TH
Một vật trượt trên mặt phẳng ngang. Áp lực lên mặt là 30 $N$, hệ số ma sát trượt 0,4. Tính lực ma sát trượt.
\choice
{12 $N$}
{13 $N$}
{\True 12 $N$}
{24 $N$}
\loigiai{
Chọn C. F_mst=μN=0.4·30=12 $N$, ngược chiều chuyển động.
}
\end{ex}

% ===== Câu 8 =====
% ID: AIQ_L10C3_FULL_0818
% Mức: TH
\begin{ex}
% ID: AIQ_L10C3_FULL_0818
% Mức: TH
Một vật trượt trên mặt phẳng ngang. Áp lực lên mặt là 40 $N$, hệ số ma sát trượt 0,5. Tính lực ma sát trượt.
\choice
{20 $N$}
{21 $N$}
{19 $N$}
{\True 20 $N$}
\loigiai{
Chọn D. F_mst=μN=0.5·40=20 $N$, ngược chiều chuyển động.
}
\end{ex}

% ===== Câu 9 =====
% ID: AIQ_L10C3_FULL_0819
% Mức: VD
\begin{ex}
% ID: AIQ_L10C3_FULL_0819
% Mức: VD
Một vật trượt trên mặt phẳng ngang. Áp lực lên mặt là 50 $N$, hệ số ma sát trượt 0,2. Tính lực ma sát trượt.
\choice
{\True 10 $N$}
{11 $N$}
{9 $N$}
{20 $N$}
\loigiai{
Chọn A. F_mst=μN=0.2·50=10 $N$, ngược chiều chuyển động.
}
\end{ex}

% ===== Câu 10 =====
% ID: AIQ_L10C3_FULL_0820
% Mức: VD
\begin{ex}
% ID: AIQ_L10C3_FULL_0820
% Mức: VD
Một vật trượt trên mặt phẳng ngang. Áp lực lên mặt là 60 $N$, hệ số ma sát trượt 0,3. Tính lực ma sát trượt.
\choice
{18 $N$}
{\True 18 $N$}
{17 $N$}
{36 $N$}
\loigiai{
Chọn B. F_mst=μN=0.3·60=18 $N$, ngược chiều chuyển động.
}
\end{ex}

% ===== Câu 11 =====
% ID: AIQ_L10C3_FULL_0821
% Mức: TH
\begin{ex}
% ID: AIQ_L10C3_FULL_0821
% Mức: TH
Một vật trượt trên mặt phẳng ngang. Áp lực lên mặt là 50 $N$, hệ số ma sát trượt 0,3. Tính lực ma sát trượt.
\choiceTF
{\True Giá trị cần tìm là $15\,\text{N}$.}
{Giá trị cần tìm là $16\,\text{N}$.}
{\True Phải xác định đúng các lực tác dụng và chọn chiều dương trước khi lập phương trình.}
{Có thể bỏ qua điều kiện áp dụng và chiều của lực mà kết quả vẫn luôn đúng.}
\loigiai{
\begin{itemchoice}
\item (Đúng) Giá trị cần tìm là $15\,\text{N}$. (Vì): F_mst=μN=0.3·50=15 $N$, ngược chiều chuyển động. b) Sai vì sai giá trị. c) Đúng. d) Sai vì phải kiểm tra điều kiện áp dụng và chiều lực
\item (Sai) Giá trị cần tìm là $16\,\text{N}$.
\item (Đúng) Phải xác định đúng các lực tác dụng và chọn chiều dương trước khi lập phương trình.
\item (Sai) Có thể bỏ qua điều kiện áp dụng và chiều của lực mà kết quả vẫn luôn đúng.
\end{itemchoice}
}
\end{ex}

% ===== Câu 12 =====
% ID: AIQ_L10C3_FULL_0822
% Mức: TH
\begin{ex}
% ID: AIQ_L10C3_FULL_0822
% Mức: TH
Một vật trượt trên mặt phẳng ngang. Áp lực lên mặt là 60 $N$, hệ số ma sát trượt 0,4. Tính lực ma sát trượt.
\choiceTF
{\True Giá trị cần tìm là $24\,\text{N}$.}
{Giá trị cần tìm là $25\,\text{N}$.}
{\True Phải xác định đúng các lực tác dụng và chọn chiều dương trước khi lập phương trình.}
{Có thể bỏ qua điều kiện áp dụng và chiều của lực mà kết quả vẫn luôn đúng.}
\loigiai{
\begin{itemchoice}
\item (Đúng) Giá trị cần tìm là $24\,\text{N}$. (Vì): F_mst=μN=0.4·60=24 $N$, ngược chiều chuyển động. b) Sai vì sai giá trị. c) Đúng. d) Sai vì phải kiểm tra điều kiện áp dụng và chiều lực
\item (Sai) Giá trị cần tìm là $25\,\text{N}$.
\item (Đúng) Phải xác định đúng các lực tác dụng và chọn chiều dương trước khi lập phương trình.
\item (Sai) Có thể bỏ qua điều kiện áp dụng và chiều của lực mà kết quả vẫn luôn đúng.
\end{itemchoice}
}
\end{ex}

% ===== Câu 13 =====
% ID: AIQ_L10C3_FULL_0823
% Mức: VD
\begin{ex}
% ID: AIQ_L10C3_FULL_0823
% Mức: VD
Một vật trượt trên mặt phẳng ngang. Áp lực lên mặt là 20 $N$, hệ số ma sát trượt 0,5. Tính lực ma sát trượt.
\choiceTF
{\True Giá trị cần tìm là $10\,\text{N}$.}
{Giá trị cần tìm là $11\,\text{N}$.}
{\True Phải xác định đúng các lực tác dụng và chọn chiều dương trước khi lập phương trình.}
{Có thể bỏ qua điều kiện áp dụng và chiều của lực mà kết quả vẫn luôn đúng.}
\loigiai{
\begin{itemchoice}
\item (Đúng) Giá trị cần tìm là $10\,\text{N}$. (Vì): F_mst=μN=0.5·20=10 $N$, ngược chiều chuyển động. b) Sai vì sai giá trị. c) Đúng. d) Sai vì phải kiểm tra điều kiện áp dụng và chiều lực
\item (Sai) Giá trị cần tìm là $11\,\text{N}$.
\item (Đúng) Phải xác định đúng các lực tác dụng và chọn chiều dương trước khi lập phương trình.
\item (Sai) Có thể bỏ qua điều kiện áp dụng và chiều của lực mà kết quả vẫn luôn đúng.
\end{itemchoice}
}
\end{ex}

% ===== Câu 14 =====
% ID: AIQ_L10C3_FULL_0824
% Mức: VD
\begin{ex}
% ID: AIQ_L10C3_FULL_0824
% Mức: VD
Một vật trượt trên mặt phẳng ngang. Áp lực lên mặt là 30 $N$, hệ số ma sát trượt 0,2. Tính lực ma sát trượt.
\choiceTF
{\True Giá trị cần tìm là $6\,\text{N}$.}
{Giá trị cần tìm là $7\,\text{N}$.}
{\True Phải xác định đúng các lực tác dụng và chọn chiều dương trước khi lập phương trình.}
{Có thể bỏ qua điều kiện áp dụng và chiều của lực mà kết quả vẫn luôn đúng.}
\loigiai{
\begin{itemchoice}
\item (Đúng) Giá trị cần tìm là $6\,\text{N}$. (Vì): F_mst=μN=0.2·30=6 $N$, ngược chiều chuyển động. b) Sai vì sai giá trị. c) Đúng. d) Sai vì phải kiểm tra điều kiện áp dụng và chiều lực
\item (Sai) Giá trị cần tìm là $7\,\text{N}$.
\item (Đúng) Phải xác định đúng các lực tác dụng và chọn chiều dương trước khi lập phương trình.
\item (Sai) Có thể bỏ qua điều kiện áp dụng và chiều của lực mà kết quả vẫn luôn đúng.
\end{itemchoice}
}
\end{ex}

% ===== Câu 15 =====
% ID: AIQ_L10C3_FULL_0825
% Mức: VD
\begin{ex}
% ID: AIQ_L10C3_FULL_0825
% Mức: VD
Một vật trượt trên mặt phẳng ngang. Áp lực lên mặt là 40 $N$, hệ số ma sát trượt 0,3. Tính lực ma sát trượt.
\choiceTF
{\True Giá trị cần tìm là $12\,\text{N}$.}
{Giá trị cần tìm là $13\,\text{N}$.}
{\True Phải xác định đúng các lực tác dụng và chọn chiều dương trước khi lập phương trình.}
{Có thể bỏ qua điều kiện áp dụng và chiều của lực mà kết quả vẫn luôn đúng.}
\loigiai{
\begin{itemchoice}
\item (Đúng) Giá trị cần tìm là $12\,\text{N}$. (Vì): F_mst=μN=0.3·40=12 $N$, ngược chiều chuyển động. b) Sai vì sai giá trị. c) Đúng. d) Sai vì phải kiểm tra điều kiện áp dụng và chiều lực
\item (Sai) Giá trị cần tìm là $13\,\text{N}$.
\item (Đúng) Phải xác định đúng các lực tác dụng và chọn chiều dương trước khi lập phương trình.
\item (Sai) Có thể bỏ qua điều kiện áp dụng và chiều của lực mà kết quả vẫn luôn đúng.
\end{itemchoice}
}
\end{ex}

% ===== Câu 16 =====
% ID: AIQ_L10C3_FULL_0826
% Mức: TH
\begin{ex}
% ID: AIQ_L10C3_FULL_0826
% Mức: TH
Một vật trượt trên mặt phẳng ngang. Áp lực lên mặt là 30 $N$, hệ số ma sát trượt 0,4. Tính lực ma sát trượt.
\shortans{12}
\loigiai{
F_mst=μN=0.4·30=12 $N$, ngược chiều chuyển động.
}
\end{ex}

% ===== Câu 17 =====
% ID: AIQ_L10C3_FULL_0827
% Mức: TH
\begin{ex}
% ID: AIQ_L10C3_FULL_0827
% Mức: TH
Một vật trượt trên mặt phẳng ngang. Áp lực lên mặt là 40 $N$, hệ số ma sát trượt 0,5. Tính lực ma sát trượt.
\shortans{20}
\loigiai{
F_mst=μN=0.5·40=20 $N$, ngược chiều chuyển động.
}
\end{ex}

% ===== Câu 18 =====
% ID: AIQ_L10C3_FULL_0828
% Mức: VD
\begin{ex}
% ID: AIQ_L10C3_FULL_0828
% Mức: VD
Một vật trượt trên mặt phẳng ngang. Áp lực lên mặt là 50 $N$, hệ số ma sát trượt 0,2. Tính lực ma sát trượt.
\shortans{10}
\loigiai{
F_mst=μN=0.2·50=10 $N$, ngược chiều chuyển động.
}
\end{ex}

% ===== Câu 19 =====
% ID: AIQ_L10C3_FULL_0829
% Mức: VD
\begin{ex}
% ID: AIQ_L10C3_FULL_0829
% Mức: VD
Một vật trượt trên mặt phẳng ngang. Áp lực lên mặt là 60 $N$, hệ số ma sát trượt 0,3. Tính lực ma sát trượt.
\shortans{18}
\loigiai{
F_mst=μN=0.3·60=18 $N$, ngược chiều chuyển động.
}
\end{ex}

% ===== Câu 20 =====
% ID: AIQ_L10C3_FULL_0830
% Mức: VD
\begin{ex}
% ID: AIQ_L10C3_FULL_0830
% Mức: VD
Một vật trượt trên mặt phẳng ngang. Áp lực lên mặt là 20 $N$, hệ số ma sát trượt 0,4. Tính lực ma sát trượt.
\shortans{8}
\loigiai{
F_mst=μN=0.4·20=8 $N$, ngược chiều chuyển động.
}
\end{ex}

% ===== Câu 21 =====
% ID: AIQ_L10C3_FULL_0831
% Mức: VDC
\begin{ex}
% ID: AIQ_L10C3_FULL_0831
% Mức: VDC
Một vật trượt trên mặt phẳng ngang. Áp lực lên mặt là 30 $N$, hệ số ma sát trượt 0,5. Tính lực ma sát trượt.
\shortans{15}
\loigiai{
F_mst=μN=0.5·30=15 $N$, ngược chiều chuyển động.
}
\end{ex}

% ===== Câu 22 =====
% ID: AIQ_L10C3_FULL_0832
% Mức: VD
\begin{bt}
% ID: AIQ_L10C3_FULL_0832
% Mức: VD
Một vật trượt trên mặt phẳng ngang. Áp lực lên mặt là 60 $N$, hệ số ma sát trượt 0,5. Tính lực ma sát trượt.
\loigiai{
Phân tích lực, chọn hệ quy chiếu và áp dụng định luật hoặc quy tắc phù hợp. F_mst=μN=0.5·60=30 $N$, ngược chiều chuyển động.
}
\end{bt}

% ===== Câu 23 =====
% ID: AIQ_L10C3_FULL_0833
% Mức: VD
\begin{bt}
% ID: AIQ_L10C3_FULL_0833
% Mức: VD
Một vật trượt trên mặt phẳng ngang. Áp lực lên mặt là 20 $N$, hệ số ma sát trượt 0,2. Tính lực ma sát trượt.
\loigiai{
Phân tích lực, chọn hệ quy chiếu và áp dụng định luật hoặc quy tắc phù hợp. F_mst=μN=0.2·20=4 $N$, ngược chiều chuyển động.
}
\end{bt}

% ===== Câu 24 =====
% ID: AIQ_L10C3_FULL_0834
% Mức: VDC
\begin{bt}
% ID: AIQ_L10C3_FULL_0834
% Mức: VDC
Một vật trượt trên mặt phẳng ngang. Áp lực lên mặt là 30 $N$, hệ số ma sát trượt 0,3. Tính lực ma sát trượt.
\loigiai{
Phân tích lực, chọn hệ quy chiếu và áp dụng định luật hoặc quy tắc phù hợp. F_mst=μN=0.3·30=9 $N$, ngược chiều chuyển động.
}
\end{bt}

% ===== Câu 25 =====
% ID: AIQ_L10C3_FULL_0835
% Mức: VDC
\begin{bt}
% ID: AIQ_L10C3_FULL_0835
% Mức: VDC
Một vật trượt trên mặt phẳng ngang. Áp lực lên mặt là 40 $N$, hệ số ma sát trượt 0,4. Tính lực ma sát trượt.
\loigiai{
Phân tích lực, chọn hệ quy chiếu và áp dụng định luật hoặc quy tắc phù hợp. F_mst=μN=0.4·40=16 $N$, ngược chiều chuyển động.
}
\end{bt}

% ===== Câu 26 =====
% ID: AIQ_L10C3_FULL_0836
% Mức: VDC
\begin{bt}
% ID: AIQ_L10C3_FULL_0836
% Mức: VDC
Một vật trượt trên mặt phẳng ngang. Áp lực lên mặt là 50 $N$, hệ số ma sát trượt 0,5. Tính lực ma sát trượt.
\loigiai{
Phân tích lực, chọn hệ quy chiếu và áp dụng định luật hoặc quy tắc phù hợp. F_mst=μN=0.5·50=25 $N$, ngược chiều chuyển động.
}
\end{bt}

% ===== Câu 27 =====
% ID: AIQ_L10C3_FULL_0837
% Mức: VDC
\begin{bt}
% ID: AIQ_L10C3_FULL_0837
% Mức: VDC
Một vật trượt trên mặt phẳng ngang. Áp lực lên mặt là 60 $N$, hệ số ma sát trượt 0,2. Tính lực ma sát trượt.
\loigiai{
Phân tích lực, chọn hệ quy chiếu và áp dụng định luật hoặc quy tắc phù hợp. F_mst=μN=0.2·60=12 $N$, ngược chiều chuyển động.
}
\end{bt}

% ===== Câu 28 =====
% ID: AIQ_L10C3_FULL_0838
% Mức: NB
\dangbt{Tính lực ma sát trượt và hệ số ma sát}
\begin{ex}
% ID: AIQ_L10C3_FULL_0838
% Mức: NB
Lực ma sát trượt giữa vật và sàn là 4 $N$, áp lực là $20\,\text{N}$. Tính hệ số ma sát trượt.
\choice
{0,2}
{\True 0,2}
{0}
{0,4}
\loigiai{
Chọn B. μ=F_mst/N=4/20=0,2.
}
\end{ex}

% ===== Câu 29 =====
% ID: AIQ_L10C3_FULL_0839
% Mức: NB
\begin{ex}
% ID: AIQ_L10C3_FULL_0839
% Mức: NB
Lực ma sát trượt giữa vật và sàn là 6 $N$, áp lực là $30\,\text{N}$. Tính hệ số ma sát trượt.
\choice
{0,2}
{1,2}
{\True 0,2}
{0,4}
\loigiai{
Chọn C. μ=F_mst/N=6/30=0,2.
}
\end{ex}

% ===== Câu 30 =====
% ID: AIQ_L10C3_FULL_0840
% Mức: NB
\begin{ex}
% ID: AIQ_L10C3_FULL_0840
% Mức: NB
Lực ma sát trượt giữa vật và sàn là 8 $N$, áp lực là $40\,\text{N}$. Tính hệ số ma sát trượt.
\choice
{0,2}
{1,2}
{0}
{\True 0,2}
\loigiai{
Chọn D. μ=F_mst/N=8/40=0,2.
}
\end{ex}

% ===== Câu 31 =====
% ID: AIQ_L10C3_FULL_0841
% Mức: NB
\begin{ex}
% ID: AIQ_L10C3_FULL_0841
% Mức: NB
Lực ma sát trượt giữa vật và sàn là 10 $N$, áp lực là $50\,\text{N}$. Tính hệ số ma sát trượt.
\choice
{\True 0,2}
{1,2}
{0}
{0,4}
\loigiai{
Chọn A. μ=F_mst/N=10/50=0,2.
}
\end{ex}

% ===== Câu 32 =====
% ID: AIQ_L10C3_FULL_0842
% Mức: TH
\begin{ex}
% ID: AIQ_L10C3_FULL_0842
% Mức: TH
Lực ma sát trượt giữa vật và sàn là 12 $N$, áp lực là $60\,\text{N}$. Tính hệ số ma sát trượt.
\choice
{0,2}
{\True 0,2}
{0}
{0,4}
\loigiai{
Chọn B. μ=F_mst/N=12/60=0,2.
}
\end{ex}

% ===== Câu 33 =====
% ID: AIQ_L10C3_FULL_0843
% Mức: TH
\begin{ex}
% ID: AIQ_L10C3_FULL_0843
% Mức: TH
Lực ma sát trượt giữa vật và sàn là 4 $N$, áp lực là $20\,\text{N}$. Tính hệ số ma sát trượt.
\choice
{0,2}
{1,2}
{\True 0,2}
{0,4}
\loigiai{
Chọn C. μ=F_mst/N=4/20=0,2.
}
\end{ex}

% ===== Câu 34 =====
% ID: AIQ_L10C3_FULL_0844
% Mức: TH
\begin{ex}
% ID: AIQ_L10C3_FULL_0844
% Mức: TH
Lực ma sát trượt giữa vật và sàn là 6 $N$, áp lực là $30\,\text{N}$. Tính hệ số ma sát trượt.
\choice
{0,2}
{1,2}
{0}
{\True 0,2}
\loigiai{
Chọn D. μ=F_mst/N=6/30=0,2.
}
\end{ex}

% ===== Câu 35 =====
% ID: AIQ_L10C3_FULL_0845
% Mức: TH
\begin{ex}
% ID: AIQ_L10C3_FULL_0845
% Mức: TH
Lực ma sát trượt giữa vật và sàn là 8 $N$, áp lực là $40\,\text{N}$. Tính hệ số ma sát trượt.
\choice
{\True 0,2}
{1,2}
{0}
{0,4}
\loigiai{
Chọn A. μ=F_mst/N=8/40=0,2.
}
\end{ex}

% ===== Câu 36 =====
% ID: AIQ_L10C3_FULL_0846
% Mức: VD
\begin{ex}
% ID: AIQ_L10C3_FULL_0846
% Mức: VD
Lực ma sát trượt giữa vật và sàn là 10 $N$, áp lực là $50\,\text{N}$. Tính hệ số ma sát trượt.
\choice
{0,2}
{\True 0,2}
{0}
{0,4}
\loigiai{
Chọn B. μ=F_mst/N=10/50=0,2.
}
\end{ex}

% ===== Câu 37 =====
% ID: AIQ_L10C3_FULL_0847
% Mức: VD
\begin{ex}
% ID: AIQ_L10C3_FULL_0847
% Mức: VD
Lực ma sát trượt giữa vật và sàn là 12 $N$, áp lực là $60\,\text{N}$. Tính hệ số ma sát trượt.
\choice
{0,2}
{1,2}
{\True 0,2}
{0,4}
\loigiai{
Chọn C. μ=F_mst/N=12/60=0,2.
}
\end{ex}

% ===== Câu 38 =====
% ID: AIQ_L10C3_FULL_0848
% Mức: TH
\begin{ex}
% ID: AIQ_L10C3_FULL_0848
% Mức: TH
Lực ma sát trượt giữa vật và sàn là 10 $N$, áp lực là $50\,\text{N}$. Tính hệ số ma sát trượt.
\choiceTF
{\True Giá trị cần tìm là 0,2.}
{Giá trị cần tìm là 1,2.}
{\True Phải xác định đúng các lực tác dụng và chọn chiều dương trước khi lập phương trình.}
{Có thể bỏ qua điều kiện áp dụng và chiều của lực mà kết quả vẫn luôn đúng.}
\loigiai{
\begin{itemchoice}
\item (Đúng) Giá trị cần tìm là 0,2. (Vì): μ=F_mst/N=10/50=0,2. b) Sai vì sai giá trị. c) Đúng. d) Sai vì phải kiểm tra điều kiện áp dụng và chiều lực
\item (Sai) Giá trị cần tìm là 1,2.
\item (Đúng) Phải xác định đúng các lực tác dụng và chọn chiều dương trước khi lập phương trình.
\item (Sai) Có thể bỏ qua điều kiện áp dụng và chiều của lực mà kết quả vẫn luôn đúng.
\end{itemchoice}
}
\end{ex}

% ===== Câu 39 =====
% ID: AIQ_L10C3_FULL_0849
% Mức: TH
\begin{ex}
% ID: AIQ_L10C3_FULL_0849
% Mức: TH
Lực ma sát trượt giữa vật và sàn là 12 $N$, áp lực là $60\,\text{N}$. Tính hệ số ma sát trượt.
\choiceTF
{\True Giá trị cần tìm là 0,2.}
{Giá trị cần tìm là 1,2.}
{\True Phải xác định đúng các lực tác dụng và chọn chiều dương trước khi lập phương trình.}
{Có thể bỏ qua điều kiện áp dụng và chiều của lực mà kết quả vẫn luôn đúng.}
\loigiai{
\begin{itemchoice}
\item (Đúng) Giá trị cần tìm là 0,2. (Vì): μ=F_mst/N=12/60=0,2. b) Sai vì sai giá trị. c) Đúng. d) Sai vì phải kiểm tra điều kiện áp dụng và chiều lực
\item (Sai) Giá trị cần tìm là 1,2.
\item (Đúng) Phải xác định đúng các lực tác dụng và chọn chiều dương trước khi lập phương trình.
\item (Sai) Có thể bỏ qua điều kiện áp dụng và chiều của lực mà kết quả vẫn luôn đúng.
\end{itemchoice}
}
\end{ex}

% ===== Câu 40 =====
% ID: AIQ_L10C3_FULL_0850
% Mức: VD
\begin{ex}
% ID: AIQ_L10C3_FULL_0850
% Mức: VD
Lực ma sát trượt giữa vật và sàn là 4 $N$, áp lực là $20\,\text{N}$. Tính hệ số ma sát trượt.
\choiceTF
{\True Giá trị cần tìm là 0,2.}
{Giá trị cần tìm là 1,2.}
{\True Phải xác định đúng các lực tác dụng và chọn chiều dương trước khi lập phương trình.}
{Có thể bỏ qua điều kiện áp dụng và chiều của lực mà kết quả vẫn luôn đúng.}
\loigiai{
\begin{itemchoice}
\item (Đúng) Giá trị cần tìm là 0,2. (Vì): μ=F_mst/N=4/20=0,2. b) Sai vì sai giá trị. c) Đúng. d) Sai vì phải kiểm tra điều kiện áp dụng và chiều lực
\item (Sai) Giá trị cần tìm là 1,2.
\item (Đúng) Phải xác định đúng các lực tác dụng và chọn chiều dương trước khi lập phương trình.
\item (Sai) Có thể bỏ qua điều kiện áp dụng và chiều của lực mà kết quả vẫn luôn đúng.
\end{itemchoice}
}
\end{ex}

% ===== Câu 41 =====
% ID: AIQ_L10C3_FULL_0851
% Mức: VD
\begin{ex}
% ID: AIQ_L10C3_FULL_0851
% Mức: VD
Lực ma sát trượt giữa vật và sàn là 6 $N$, áp lực là $30\,\text{N}$. Tính hệ số ma sát trượt.
\choiceTF
{\True Giá trị cần tìm là 0,2.}
{Giá trị cần tìm là 1,2.}
{\True Phải xác định đúng các lực tác dụng và chọn chiều dương trước khi lập phương trình.}
{Có thể bỏ qua điều kiện áp dụng và chiều của lực mà kết quả vẫn luôn đúng.}
\loigiai{
\begin{itemchoice}
\item (Đúng) Giá trị cần tìm là 0,2. (Vì): μ=F_mst/N=6/30=0,2. b) Sai vì sai giá trị. c) Đúng. d) Sai vì phải kiểm tra điều kiện áp dụng và chiều lực
\item (Sai) Giá trị cần tìm là 1,2.
\item (Đúng) Phải xác định đúng các lực tác dụng và chọn chiều dương trước khi lập phương trình.
\item (Sai) Có thể bỏ qua điều kiện áp dụng và chiều của lực mà kết quả vẫn luôn đúng.
\end{itemchoice}
}
\end{ex}

% ===== Câu 42 =====
% ID: AIQ_L10C3_FULL_0852
% Mức: VD
\begin{ex}
% ID: AIQ_L10C3_FULL_0852
% Mức: VD
Lực ma sát trượt giữa vật và sàn là 8 $N$, áp lực là $40\,\text{N}$. Tính hệ số ma sát trượt.
\choiceTF
{\True Giá trị cần tìm là 0,2.}
{Giá trị cần tìm là 1,2.}
{\True Phải xác định đúng các lực tác dụng và chọn chiều dương trước khi lập phương trình.}
{Có thể bỏ qua điều kiện áp dụng và chiều của lực mà kết quả vẫn luôn đúng.}
\loigiai{
\begin{itemchoice}
\item (Đúng) Giá trị cần tìm là 0,2. (Vì): μ=F_mst/N=8/40=0,2. b) Sai vì sai giá trị. c) Đúng. d) Sai vì phải kiểm tra điều kiện áp dụng và chiều lực
\item (Sai) Giá trị cần tìm là 1,2.
\item (Đúng) Phải xác định đúng các lực tác dụng và chọn chiều dương trước khi lập phương trình.
\item (Sai) Có thể bỏ qua điều kiện áp dụng và chiều của lực mà kết quả vẫn luôn đúng.
\end{itemchoice}
}
\end{ex}

% ===== Câu 43 =====
% ID: AIQ_L10C3_FULL_0853
% Mức: TH
\begin{ex}
% ID: AIQ_L10C3_FULL_0853
% Mức: TH
Lực ma sát trượt giữa vật và sàn là 6 $N$, áp lực là $30\,\text{N}$. Tính hệ số ma sát trượt.
\shortans{0,2}
\loigiai{
μ=F_mst/N=6/30=0,2.
}
\end{ex}

% ===== Câu 44 =====
% ID: AIQ_L10C3_FULL_0854
% Mức: TH
\begin{ex}
% ID: AIQ_L10C3_FULL_0854
% Mức: TH
Lực ma sát trượt giữa vật và sàn là 8 $N$, áp lực là $40\,\text{N}$. Tính hệ số ma sát trượt.
\shortans{0,2}
\loigiai{
μ=F_mst/N=8/40=0,2.
}
\end{ex}

% ===== Câu 45 =====
% ID: AIQ_L10C3_FULL_0855
% Mức: VD
\begin{ex}
% ID: AIQ_L10C3_FULL_0855
% Mức: VD
Lực ma sát trượt giữa vật và sàn là 10 $N$, áp lực là $50\,\text{N}$. Tính hệ số ma sát trượt.
\shortans{0,2}
\loigiai{
μ=F_mst/N=10/50=0,2.
}
\end{ex}

% ===== Câu 46 =====
% ID: AIQ_L10C3_FULL_0856
% Mức: VD
\begin{ex}
% ID: AIQ_L10C3_FULL_0856
% Mức: VD
Lực ma sát trượt giữa vật và sàn là 12 $N$, áp lực là $60\,\text{N}$. Tính hệ số ma sát trượt.
\shortans{0,2}
\loigiai{
μ=F_mst/N=12/60=0,2.
}
\end{ex}

% ===== Câu 47 =====
% ID: AIQ_L10C3_FULL_0857
% Mức: VD
\begin{ex}
% ID: AIQ_L10C3_FULL_0857
% Mức: VD
Lực ma sát trượt giữa vật và sàn là 4 $N$, áp lực là $20\,\text{N}$. Tính hệ số ma sát trượt.
\shortans{0,2}
\loigiai{
μ=F_mst/N=4/20=0,2.
}
\end{ex}

% ===== Câu 48 =====
% ID: AIQ_L10C3_FULL_0858
% Mức: VDC
\begin{ex}
% ID: AIQ_L10C3_FULL_0858
% Mức: VDC
Lực ma sát trượt giữa vật và sàn là 6 $N$, áp lực là $30\,\text{N}$. Tính hệ số ma sát trượt.
\shortans{0,2}
\loigiai{
μ=F_mst/N=6/30=0,2.
}
\end{ex}

% ===== Câu 49 =====
% ID: AIQ_L10C3_FULL_0859
% Mức: VD
\begin{bt}
% ID: AIQ_L10C3_FULL_0859
% Mức: VD
Lực ma sát trượt giữa vật và sàn là 12 $N$, áp lực là $60\,\text{N}$. Tính hệ số ma sát trượt.
\loigiai{
Phân tích lực, chọn hệ quy chiếu và áp dụng định luật hoặc quy tắc phù hợp. μ=F_mst/N=12/60=0,2.
}
\end{bt}

% ===== Câu 50 =====
% ID: AIQ_L10C3_FULL_0860
% Mức: VD
\begin{bt}
% ID: AIQ_L10C3_FULL_0860
% Mức: VD
Lực ma sát trượt giữa vật và sàn là 4 $N$, áp lực là $20\,\text{N}$. Tính hệ số ma sát trượt.
\loigiai{
Phân tích lực, chọn hệ quy chiếu và áp dụng định luật hoặc quy tắc phù hợp. μ=F_mst/N=4/20=0,2.
}
\end{bt}

% ===== Câu 51 =====
% ID: AIQ_L10C3_FULL_0861
% Mức: VDC
\begin{bt}
% ID: AIQ_L10C3_FULL_0861
% Mức: VDC
Lực ma sát trượt giữa vật và sàn là 6 $N$, áp lực là $30\,\text{N}$. Tính hệ số ma sát trượt.
\loigiai{
Phân tích lực, chọn hệ quy chiếu và áp dụng định luật hoặc quy tắc phù hợp. μ=F_mst/N=6/30=0,2.
}
\end{bt}

% ===== Câu 52 =====
% ID: AIQ_L10C3_FULL_0862
% Mức: VDC
\begin{bt}
% ID: AIQ_L10C3_FULL_0862
% Mức: VDC
Lực ma sát trượt giữa vật và sàn là 8 $N$, áp lực là $40\,\text{N}$. Tính hệ số ma sát trượt.
\loigiai{
Phân tích lực, chọn hệ quy chiếu và áp dụng định luật hoặc quy tắc phù hợp. μ=F_mst/N=8/40=0,2.
}
\end{bt}

% ===== Câu 53 =====
% ID: AIQ_L10C3_FULL_0863
% Mức: VDC
\begin{bt}
% ID: AIQ_L10C3_FULL_0863
% Mức: VDC
Lực ma sát trượt giữa vật và sàn là 10 $N$, áp lực là $50\,\text{N}$. Tính hệ số ma sát trượt.
\loigiai{
Phân tích lực, chọn hệ quy chiếu và áp dụng định luật hoặc quy tắc phù hợp. μ=F_mst/N=10/50=0,2.
}
\end{bt}

% ===== Câu 54 =====
% ID: AIQ_L10C3_FULL_0864
% Mức: VDC
\begin{bt}
% ID: AIQ_L10C3_FULL_0864
% Mức: VDC
Lực ma sát trượt giữa vật và sàn là 12 $N$, áp lực là $60\,\text{N}$. Tính hệ số ma sát trượt.
\loigiai{
Phân tích lực, chọn hệ quy chiếu và áp dụng định luật hoặc quy tắc phù hợp. μ=F_mst/N=12/60=0,2.
}
\end{bt}

% ===== Câu 55 =====
% ID: AIQ_L10C3_FULL_0865
% Mức: NB
\dangbt{Xác định lực ma sát nghỉ và điều kiện trượt}
\begin{ex}
% ID: AIQ_L10C3_FULL_0865
% Mức: NB
Vật đứng yên trên sàn ngang dù bị kéo ngang lực 5 $N$; lực kéo chưa vượt ma sát nghỉ cực đại. Tính độ lớn lực ma sát nghỉ.
\choice
{5 $N$}
{6 $N$}
{\True 5 $N$}
{10 $N$}
\loigiai{
Chọn C. Khi vật còn đứng yên, ma sát nghỉ tự điều chỉnh cân bằng lực kéo nên F_msn= $5\,\text{N}$.
}
\end{ex}

% ===== Câu 56 =====
% ID: AIQ_L10C3_FULL_0866
% Mức: NB
\begin{ex}
% ID: AIQ_L10C3_FULL_0866
% Mức: NB
Vật đứng yên trên sàn ngang dù bị kéo ngang lực 6 $N$; lực kéo chưa vượt ma sát nghỉ cực đại. Tính độ lớn lực ma sát nghỉ.
\choice
{6 $N$}
{7 $N$}
{5 $N$}
{\True 6 $N$}
\loigiai{
Chọn D. Khi vật còn đứng yên, ma sát nghỉ tự điều chỉnh cân bằng lực kéo nên F_msn= $6\,\text{N}$.
}
\end{ex}

% ===== Câu 57 =====
% ID: AIQ_L10C3_FULL_0867
% Mức: NB
\begin{ex}
% ID: AIQ_L10C3_FULL_0867
% Mức: NB
Vật đứng yên trên sàn ngang dù bị kéo ngang lực 7 $N$; lực kéo chưa vượt ma sát nghỉ cực đại. Tính độ lớn lực ma sát nghỉ.
\choice
{\True 7 $N$}
{8 $N$}
{6 $N$}
{14 $N$}
\loigiai{
Chọn A. Khi vật còn đứng yên, ma sát nghỉ tự điều chỉnh cân bằng lực kéo nên F_msn= $7\,\text{N}$.
}
\end{ex}

% ===== Câu 58 =====
% ID: AIQ_L10C3_FULL_0868
% Mức: NB
\begin{ex}
% ID: AIQ_L10C3_FULL_0868
% Mức: NB
Vật đứng yên trên sàn ngang dù bị kéo ngang lực 8 $N$; lực kéo chưa vượt ma sát nghỉ cực đại. Tính độ lớn lực ma sát nghỉ.
\choice
{8 $N$}
{\True 8 $N$}
{7 $N$}
{16 $N$}
\loigiai{
Chọn B. Khi vật còn đứng yên, ma sát nghỉ tự điều chỉnh cân bằng lực kéo nên F_msn= $8\,\text{N}$.
}
\end{ex}

% ===== Câu 59 =====
% ID: AIQ_L10C3_FULL_0869
% Mức: TH
\begin{ex}
% ID: AIQ_L10C3_FULL_0869
% Mức: TH
Vật đứng yên trên sàn ngang dù bị kéo ngang lực 9 $N$; lực kéo chưa vượt ma sát nghỉ cực đại. Tính độ lớn lực ma sát nghỉ.
\choice
{9 $N$}
{10 $N$}
{\True 9 $N$}
{18 $N$}
\loigiai{
Chọn C. Khi vật còn đứng yên, ma sát nghỉ tự điều chỉnh cân bằng lực kéo nên F_msn= $9\,\text{N}$.
}
\end{ex}

% ===== Câu 60 =====
% ID: AIQ_L10C3_FULL_0870
% Mức: TH
\begin{ex}
% ID: AIQ_L10C3_FULL_0870
% Mức: TH
Vật đứng yên trên sàn ngang dù bị kéo ngang lực 10 $N$; lực kéo chưa vượt ma sát nghỉ cực đại. Tính độ lớn lực ma sát nghỉ.
\choice
{10 $N$}
{11 $N$}
{9 $N$}
{\True 10 $N$}
\loigiai{
Chọn D. Khi vật còn đứng yên, ma sát nghỉ tự điều chỉnh cân bằng lực kéo nên F_msn= $10\,\text{N}$.
}
\end{ex}

% ===== Câu 61 =====
% ID: AIQ_L10C3_FULL_0871
% Mức: TH
\begin{ex}
% ID: AIQ_L10C3_FULL_0871
% Mức: TH
Vật đứng yên trên sàn ngang dù bị kéo ngang lực 11 $N$; lực kéo chưa vượt ma sát nghỉ cực đại. Tính độ lớn lực ma sát nghỉ.
\choice
{\True 11 $N$}
{12 $N$}
{10 $N$}
{22 $N$}
\loigiai{
Chọn A. Khi vật còn đứng yên, ma sát nghỉ tự điều chỉnh cân bằng lực kéo nên F_msn= $11\,\text{N}$.
}
\end{ex}

% ===== Câu 62 =====
% ID: AIQ_L10C3_FULL_0872
% Mức: TH
\begin{ex}
% ID: AIQ_L10C3_FULL_0872
% Mức: TH
Vật đứng yên trên sàn ngang dù bị kéo ngang lực 12 $N$; lực kéo chưa vượt ma sát nghỉ cực đại. Tính độ lớn lực ma sát nghỉ.
\choice
{12 $N$}
{\True 12 $N$}
{11 $N$}
{24 $N$}
\loigiai{
Chọn B. Khi vật còn đứng yên, ma sát nghỉ tự điều chỉnh cân bằng lực kéo nên F_msn= $12\,\text{N}$.
}
\end{ex}

% ===== Câu 63 =====
% ID: AIQ_L10C3_FULL_0873
% Mức: VD
\begin{ex}
% ID: AIQ_L10C3_FULL_0873
% Mức: VD
Vật đứng yên trên sàn ngang dù bị kéo ngang lực 13 $N$; lực kéo chưa vượt ma sát nghỉ cực đại. Tính độ lớn lực ma sát nghỉ.
\choice
{13 $N$}
{14 $N$}
{\True 13 $N$}
{26 $N$}
\loigiai{
Chọn C. Khi vật còn đứng yên, ma sát nghỉ tự điều chỉnh cân bằng lực kéo nên F_msn= $13\,\text{N}$.
}
\end{ex}

% ===== Câu 64 =====
% ID: AIQ_L10C3_FULL_0874
% Mức: VD
\begin{ex}
% ID: AIQ_L10C3_FULL_0874
% Mức: VD
Vật đứng yên trên sàn ngang dù bị kéo ngang lực 14 $N$; lực kéo chưa vượt ma sát nghỉ cực đại. Tính độ lớn lực ma sát nghỉ.
\choice
{14 $N$}
{15 $N$}
{13 $N$}
{\True 14 $N$}
\loigiai{
Chọn D. Khi vật còn đứng yên, ma sát nghỉ tự điều chỉnh cân bằng lực kéo nên F_msn= $14\,\text{N}$.
}
\end{ex}

% ===== Câu 65 =====
% ID: AIQ_L10C3_FULL_0875
% Mức: TH
\begin{ex}
% ID: AIQ_L10C3_FULL_0875
% Mức: TH
Vật đứng yên trên sàn ngang dù bị kéo ngang lực 8 $N$; lực kéo chưa vượt ma sát nghỉ cực đại. Tính độ lớn lực ma sát nghỉ.
\choiceTF
{\True Giá trị cần tìm là $8\,\text{N}$.}
{Giá trị cần tìm là $9\,\text{N}$.}
{\True Phải xác định đúng các lực tác dụng và chọn chiều dương trước khi lập phương trình.}
{Có thể bỏ qua điều kiện áp dụng và chiều của lực mà kết quả vẫn luôn đúng.}
\loigiai{
\begin{itemchoice}
\item (Đúng) Giá trị cần tìm là $8\,\text{N}$. (Vì): Khi vật còn đứng yên, ma sát nghỉ tự điều chỉnh cân bằng lực kéo nên F_msn= $8\,\text{N}$. b) Sai vì sai giá trị. c) Đúng. d) Sai vì phải kiểm tra điều kiện áp dụng và chiều lực
\item (Sai) Giá trị cần tìm là $9\,\text{N}$.
\item (Đúng) Phải xác định đúng các lực tác dụng và chọn chiều dương trước khi lập phương trình.
\item (Sai) Có thể bỏ qua điều kiện áp dụng và chiều của lực mà kết quả vẫn luôn đúng.
\end{itemchoice}
}
\end{ex}

% ===== Câu 66 =====
% ID: AIQ_L10C3_FULL_0876
% Mức: TH
\begin{ex}
% ID: AIQ_L10C3_FULL_0876
% Mức: TH
Vật đứng yên trên sàn ngang dù bị kéo ngang lực 9 $N$; lực kéo chưa vượt ma sát nghỉ cực đại. Tính độ lớn lực ma sát nghỉ.
\choiceTF
{\True Giá trị cần tìm là $9\,\text{N}$.}
{Giá trị cần tìm là $10\,\text{N}$.}
{\True Phải xác định đúng các lực tác dụng và chọn chiều dương trước khi lập phương trình.}
{Có thể bỏ qua điều kiện áp dụng và chiều của lực mà kết quả vẫn luôn đúng.}
\loigiai{
\begin{itemchoice}
\item (Đúng) Giá trị cần tìm là $9\,\text{N}$. (Vì): Khi vật còn đứng yên, ma sát nghỉ tự điều chỉnh cân bằng lực kéo nên F_msn= $9\,\text{N}$. b) Sai vì sai giá trị. c) Đúng. d) Sai vì phải kiểm tra điều kiện áp dụng và chiều lực
\item (Sai) Giá trị cần tìm là $10\,\text{N}$.
\item (Đúng) Phải xác định đúng các lực tác dụng và chọn chiều dương trước khi lập phương trình.
\item (Sai) Có thể bỏ qua điều kiện áp dụng và chiều của lực mà kết quả vẫn luôn đúng.
\end{itemchoice}
}
\end{ex}

% ===== Câu 67 =====
% ID: AIQ_L10C3_FULL_0877
% Mức: VD
\begin{ex}
% ID: AIQ_L10C3_FULL_0877
% Mức: VD
Vật đứng yên trên sàn ngang dù bị kéo ngang lực 10 $N$; lực kéo chưa vượt ma sát nghỉ cực đại. Tính độ lớn lực ma sát nghỉ.
\choiceTF
{\True Giá trị cần tìm là $10\,\text{N}$.}
{Giá trị cần tìm là $11\,\text{N}$.}
{\True Phải xác định đúng các lực tác dụng và chọn chiều dương trước khi lập phương trình.}
{Có thể bỏ qua điều kiện áp dụng và chiều của lực mà kết quả vẫn luôn đúng.}
\loigiai{
\begin{itemchoice}
\item (Đúng) Giá trị cần tìm là $10\,\text{N}$. (Vì): Khi vật còn đứng yên, ma sát nghỉ tự điều chỉnh cân bằng lực kéo nên F_msn= $10\,\text{N}$. b) Sai vì sai giá trị. c) Đúng. d) Sai vì phải kiểm tra điều kiện áp dụng và chiều lực
\item (Sai) Giá trị cần tìm là $11\,\text{N}$.
\item (Đúng) Phải xác định đúng các lực tác dụng và chọn chiều dương trước khi lập phương trình.
\item (Sai) Có thể bỏ qua điều kiện áp dụng và chiều của lực mà kết quả vẫn luôn đúng.
\end{itemchoice}
}
\end{ex}

% ===== Câu 68 =====
% ID: AIQ_L10C3_FULL_0878
% Mức: VD
\begin{ex}
% ID: AIQ_L10C3_FULL_0878
% Mức: VD
Vật đứng yên trên sàn ngang dù bị kéo ngang lực 11 $N$; lực kéo chưa vượt ma sát nghỉ cực đại. Tính độ lớn lực ma sát nghỉ.
\choiceTF
{\True Giá trị cần tìm là $11\,\text{N}$.}
{Giá trị cần tìm là $12\,\text{N}$.}
{\True Phải xác định đúng các lực tác dụng và chọn chiều dương trước khi lập phương trình.}
{Có thể bỏ qua điều kiện áp dụng và chiều của lực mà kết quả vẫn luôn đúng.}
\loigiai{
\begin{itemchoice}
\item (Đúng) Giá trị cần tìm là $11\,\text{N}$. (Vì): Khi vật còn đứng yên, ma sát nghỉ tự điều chỉnh cân bằng lực kéo nên F_msn= $11\,\text{N}$. b) Sai vì sai giá trị. c) Đúng. d) Sai vì phải kiểm tra điều kiện áp dụng và chiều lực
\item (Sai) Giá trị cần tìm là $12\,\text{N}$.
\item (Đúng) Phải xác định đúng các lực tác dụng và chọn chiều dương trước khi lập phương trình.
\item (Sai) Có thể bỏ qua điều kiện áp dụng và chiều của lực mà kết quả vẫn luôn đúng.
\end{itemchoice}
}
\end{ex}

% ===== Câu 69 =====
% ID: AIQ_L10C3_FULL_0879
% Mức: VD
\begin{ex}
% ID: AIQ_L10C3_FULL_0879
% Mức: VD
Vật đứng yên trên sàn ngang dù bị kéo ngang lực 12 $N$; lực kéo chưa vượt ma sát nghỉ cực đại. Tính độ lớn lực ma sát nghỉ.
\choiceTF
{\True Giá trị cần tìm là $12\,\text{N}$.}
{Giá trị cần tìm là $13\,\text{N}$.}
{\True Phải xác định đúng các lực tác dụng và chọn chiều dương trước khi lập phương trình.}
{Có thể bỏ qua điều kiện áp dụng và chiều của lực mà kết quả vẫn luôn đúng.}
\loigiai{
\begin{itemchoice}
\item (Đúng) Giá trị cần tìm là $12\,\text{N}$. (Vì): Khi vật còn đứng yên, ma sát nghỉ tự điều chỉnh cân bằng lực kéo nên F_msn= $12\,\text{N}$. b) Sai vì sai giá trị. c) Đúng. d) Sai vì phải kiểm tra điều kiện áp dụng và chiều lực
\item (Sai) Giá trị cần tìm là $13\,\text{N}$.
\item (Đúng) Phải xác định đúng các lực tác dụng và chọn chiều dương trước khi lập phương trình.
\item (Sai) Có thể bỏ qua điều kiện áp dụng và chiều của lực mà kết quả vẫn luôn đúng.
\end{itemchoice}
}
\end{ex}

% ===== Câu 70 =====
% ID: AIQ_L10C3_FULL_0880
% Mức: TH
\begin{ex}
% ID: AIQ_L10C3_FULL_0880
% Mức: TH
Vật đứng yên trên sàn ngang dù bị kéo ngang lực 11 $N$; lực kéo chưa vượt ma sát nghỉ cực đại. Tính độ lớn lực ma sát nghỉ.
\shortans{11}
\loigiai{
Khi vật còn đứng yên, ma sát nghỉ tự điều chỉnh cân bằng lực kéo nên F_msn= $11\,\text{N}$.
}
\end{ex}

% ===== Câu 71 =====
% ID: AIQ_L10C3_FULL_0881
% Mức: TH
\begin{ex}
% ID: AIQ_L10C3_FULL_0881
% Mức: TH
Vật đứng yên trên sàn ngang dù bị kéo ngang lực 12 $N$; lực kéo chưa vượt ma sát nghỉ cực đại. Tính độ lớn lực ma sát nghỉ.
\shortans{12}
\loigiai{
Khi vật còn đứng yên, ma sát nghỉ tự điều chỉnh cân bằng lực kéo nên F_msn= $12\,\text{N}$.
}
\end{ex}

% ===== Câu 72 =====
% ID: AIQ_L10C3_FULL_0882
% Mức: VD
\begin{ex}
% ID: AIQ_L10C3_FULL_0882
% Mức: VD
Vật đứng yên trên sàn ngang dù bị kéo ngang lực 13 $N$; lực kéo chưa vượt ma sát nghỉ cực đại. Tính độ lớn lực ma sát nghỉ.
\shortans{13}
\loigiai{
Khi vật còn đứng yên, ma sát nghỉ tự điều chỉnh cân bằng lực kéo nên F_msn= $13\,\text{N}$.
}
\end{ex}

% ===== Câu 73 =====
% ID: AIQ_L10C3_FULL_0883
% Mức: VD
\begin{ex}
% ID: AIQ_L10C3_FULL_0883
% Mức: VD
Vật đứng yên trên sàn ngang dù bị kéo ngang lực 14 $N$; lực kéo chưa vượt ma sát nghỉ cực đại. Tính độ lớn lực ma sát nghỉ.
\shortans{14}
\loigiai{
Khi vật còn đứng yên, ma sát nghỉ tự điều chỉnh cân bằng lực kéo nên F_msn= $14\,\text{N}$.
}
\end{ex}

% ===== Câu 74 =====
% ID: AIQ_L10C3_FULL_0884
% Mức: VD
\begin{ex}
% ID: AIQ_L10C3_FULL_0884
% Mức: VD
Vật đứng yên trên sàn ngang dù bị kéo ngang lực 5 $N$; lực kéo chưa vượt ma sát nghỉ cực đại. Tính độ lớn lực ma sát nghỉ.
\shortans{5}
\loigiai{
Khi vật còn đứng yên, ma sát nghỉ tự điều chỉnh cân bằng lực kéo nên F_msn= $5\,\text{N}$.
}
\end{ex}

% ===== Câu 75 =====
% ID: AIQ_L10C3_FULL_0885
% Mức: VDC
\begin{ex}
% ID: AIQ_L10C3_FULL_0885
% Mức: VDC
Vật đứng yên trên sàn ngang dù bị kéo ngang lực 6 $N$; lực kéo chưa vượt ma sát nghỉ cực đại. Tính độ lớn lực ma sát nghỉ.
\shortans{6}
\loigiai{
Khi vật còn đứng yên, ma sát nghỉ tự điều chỉnh cân bằng lực kéo nên F_msn= $6\,\text{N}$.
}
\end{ex}

% ===== Câu 76 =====
% ID: AIQ_L10C3_FULL_0886
% Mức: VD
\begin{bt}
% ID: AIQ_L10C3_FULL_0886
% Mức: VD
Vật đứng yên trên sàn ngang dù bị kéo ngang lực 14 $N$; lực kéo chưa vượt ma sát nghỉ cực đại. Tính độ lớn lực ma sát nghỉ.
\loigiai{
Phân tích lực, chọn hệ quy chiếu và áp dụng định luật hoặc quy tắc phù hợp. Khi vật còn đứng yên, ma sát nghỉ tự điều chỉnh cân bằng lực kéo nên F_msn= $14\,\text{N}$.
}
\end{bt}

% ===== Câu 77 =====
% ID: AIQ_L10C3_FULL_0887
% Mức: VD
\begin{bt}
% ID: AIQ_L10C3_FULL_0887
% Mức: VD
Vật đứng yên trên sàn ngang dù bị kéo ngang lực 5 $N$; lực kéo chưa vượt ma sát nghỉ cực đại. Tính độ lớn lực ma sát nghỉ.
\loigiai{
Phân tích lực, chọn hệ quy chiếu và áp dụng định luật hoặc quy tắc phù hợp. Khi vật còn đứng yên, ma sát nghỉ tự điều chỉnh cân bằng lực kéo nên F_msn= $5\,\text{N}$.
}
\end{bt}

% ===== Câu 78 =====
% ID: AIQ_L10C3_FULL_0888
% Mức: VDC
\begin{bt}
% ID: AIQ_L10C3_FULL_0888
% Mức: VDC
Vật đứng yên trên sàn ngang dù bị kéo ngang lực 6 $N$; lực kéo chưa vượt ma sát nghỉ cực đại. Tính độ lớn lực ma sát nghỉ.
\loigiai{
Phân tích lực, chọn hệ quy chiếu và áp dụng định luật hoặc quy tắc phù hợp. Khi vật còn đứng yên, ma sát nghỉ tự điều chỉnh cân bằng lực kéo nên F_msn= $6\,\text{N}$.
}
\end{bt}

% ===== Câu 79 =====
% ID: AIQ_L10C3_FULL_0889
% Mức: VDC
\begin{bt}
% ID: AIQ_L10C3_FULL_0889
% Mức: VDC
Vật đứng yên trên sàn ngang dù bị kéo ngang lực 7 $N$; lực kéo chưa vượt ma sát nghỉ cực đại. Tính độ lớn lực ma sát nghỉ.
\loigiai{
Phân tích lực, chọn hệ quy chiếu và áp dụng định luật hoặc quy tắc phù hợp. Khi vật còn đứng yên, ma sát nghỉ tự điều chỉnh cân bằng lực kéo nên F_msn= $7\,\text{N}$.
}
\end{bt}

% ===== Câu 80 =====
% ID: AIQ_L10C3_FULL_0890
% Mức: VDC
\begin{bt}
% ID: AIQ_L10C3_FULL_0890
% Mức: VDC
Vật đứng yên trên sàn ngang dù bị kéo ngang lực 8 $N$; lực kéo chưa vượt ma sát nghỉ cực đại. Tính độ lớn lực ma sát nghỉ.
\loigiai{
Phân tích lực, chọn hệ quy chiếu và áp dụng định luật hoặc quy tắc phù hợp. Khi vật còn đứng yên, ma sát nghỉ tự điều chỉnh cân bằng lực kéo nên F_msn= $8\,\text{N}$.
}
\end{bt}

% ===== Câu 81 =====
% ID: AIQ_L10C3_FULL_0891
% Mức: VDC
\begin{bt}
% ID: AIQ_L10C3_FULL_0891
% Mức: VDC
Vật đứng yên trên sàn ngang dù bị kéo ngang lực 9 $N$; lực kéo chưa vượt ma sát nghỉ cực đại. Tính độ lớn lực ma sát nghỉ.
\loigiai{
Phân tích lực, chọn hệ quy chiếu và áp dụng định luật hoặc quy tắc phù hợp. Khi vật còn đứng yên, ma sát nghỉ tự điều chỉnh cân bằng lực kéo nên F_msn= $9\,\text{N}$.
}
\end{bt}

% ===== Câu 82 =====
% ID: AIQ_L10C3_FULL_0892
% Mức: NB
\dangbt{Chuyển động trên mặt phẳng ngang có ma sát}
\begin{ex}
% ID: AIQ_L10C3_FULL_0892
% Mức: NB
Vật $2\,\text{kg}$ trên sàn ngang có μ=0,2, chịu lực kéo ngang $12\,\text{N}$. Lấy g= $10\,\text{m}$ /s^2. Tính gia tốc.
\choice
{$4 m/s^2$}
{$5 m/s^2$}
{$3 m/s^2$}
{\True $4 m/s^2$}
\loigiai{
Chọn D. F_ms=μmg=4 $N; a=(F−F_ms)/m=(12−4)/2=4\,\text{m}$ /s^2.
}
\end{ex}

% ===== Câu 83 =====
% ID: AIQ_L10C3_FULL_0893
% Mức: NB
\begin{ex}
% ID: AIQ_L10C3_FULL_0893
% Mức: NB
Vật $3\,\text{kg}$ trên sàn ngang có μ=0,3, chịu lực kéo ngang $12\,\text{N}$. Lấy g= $10\,\text{m}$ /s^2. Tính gia tốc.
\choice
{\True $1 m/s^2$}
{$2 m/s^2$}
{$0 m/s^2$}
{$2 m/s^2$}
\loigiai{
Chọn A. F_ms=μmg=9 $N; a=(F−F_ms)/m=(12−9)/3=1\,\text{m}$ /s^2.
}
\end{ex}

% ===== Câu 84 =====
% ID: AIQ_L10C3_FULL_0894
% Mức: NB
\begin{ex}
% ID: AIQ_L10C3_FULL_0894
% Mức: NB
Vật $4\,\text{kg}$ trên sàn ngang có μ=0,4, chịu lực kéo ngang $24\,\text{N}$. Lấy g= $10\,\text{m}$ /s^2. Tính gia tốc.
\choice
{$2 m/s^2$}
{\True $2 m/s^2$}
{$1 m/s^2$}
{$4 m/s^2$}
\loigiai{
Chọn B. F_ms=μmg=16 $N; a=(F−F_ms)/m=(24−16)/4=2\,\text{m}$ /s^2.
}
\end{ex}

% ===== Câu 85 =====
% ID: AIQ_L10C3_FULL_0895
% Mức: NB
\begin{ex}
% ID: AIQ_L10C3_FULL_0895
% Mức: NB
Vật $5\,\text{kg}$ trên sàn ngang có μ=0,5, chịu lực kéo ngang $40\,\text{N}$. Lấy g= $10\,\text{m}$ /s^2. Tính gia tốc.
\choice
{$3 m/s^2$}
{$4 m/s^2$}
{\True $3 m/s^2$}
{$6 m/s^2$}
\loigiai{
Chọn C. F_ms=μmg=25 $N; a=(F−F_ms)/m=(40−25)/5=3\,\text{m}$ /s^2.
}
\end{ex}

% ===== Câu 86 =====
% ID: AIQ_L10C3_FULL_0896
% Mức: TH
\begin{ex}
% ID: AIQ_L10C3_FULL_0896
% Mức: TH
Vật $6\,\text{kg}$ trên sàn ngang có μ=0,2, chịu lực kéo ngang $36\,\text{N}$. Lấy g= $10\,\text{m}$ /s^2. Tính gia tốc.
\choice
{$4 m/s^2$}
{$5 m/s^2$}
{$3 m/s^2$}
{\True $4 m/s^2$}
\loigiai{
Chọn D. F_ms=μmg=12 $N; a=(F−F_ms)/m=(36−12)/6=4\,\text{m}$ /s^2.
}
\end{ex}

% ===== Câu 87 =====
% ID: AIQ_L10C3_FULL_0897
% Mức: TH
\begin{ex}
% ID: AIQ_L10C3_FULL_0897
% Mức: TH
Vật $2\,\text{kg}$ trên sàn ngang có μ=0,3, chịu lực kéo ngang $8\,\text{N}$. Lấy g= $10\,\text{m}$ /s^2. Tính gia tốc.
\choice
{\True $1 m/s^2$}
{$2 m/s^2$}
{$0 m/s^2$}
{$2 m/s^2$}
\loigiai{
Chọn A. F_ms=μmg=6 $N; a=(F−F_ms)/m=(8−6)/2=1\,\text{m}$ /s^2.
}
\end{ex}

% ===== Câu 88 =====
% ID: AIQ_L10C3_FULL_0898
% Mức: TH
\begin{ex}
% ID: AIQ_L10C3_FULL_0898
% Mức: TH
Vật $3\,\text{kg}$ trên sàn ngang có μ=0,4, chịu lực kéo ngang $18\,\text{N}$. Lấy g= $10\,\text{m}$ /s^2. Tính gia tốc.
\choice
{$2 m/s^2$}
{\True $2 m/s^2$}
{$1 m/s^2$}
{$4 m/s^2$}
\loigiai{
Chọn B. F_ms=μmg=12 $N; a=(F−F_ms)/m=(18−12)/3=2\,\text{m}$ /s^2.
}
\end{ex}

% ===== Câu 89 =====
% ID: AIQ_L10C3_FULL_0899
% Mức: TH
\begin{ex}
% ID: AIQ_L10C3_FULL_0899
% Mức: TH
Vật $4\,\text{kg}$ trên sàn ngang có μ=0,5, chịu lực kéo ngang $32\,\text{N}$. Lấy g= $10\,\text{m}$ /s^2. Tính gia tốc.
\choice
{$3 m/s^2$}
{$4 m/s^2$}
{\True $3 m/s^2$}
{$6 m/s^2$}
\loigiai{
Chọn C. F_ms=μmg=20 $N; a=(F−F_ms)/m=(32−20)/4=3\,\text{m}$ /s^2.
}
\end{ex}

% ===== Câu 90 =====
% ID: AIQ_L10C3_FULL_0900
% Mức: VD
\begin{ex}
% ID: AIQ_L10C3_FULL_0900
% Mức: VD
Vật $5\,\text{kg}$ trên sàn ngang có μ=0,2, chịu lực kéo ngang $30\,\text{N}$. Lấy g= $10\,\text{m}$ /s^2. Tính gia tốc.
\choice
{$4 m/s^2$}
{$5 m/s^2$}
{$3 m/s^2$}
{\True $4 m/s^2$}
\loigiai{
Chọn D. F_ms=μmg=10 $N; a=(F−F_ms)/m=(30−10)/5=4\,\text{m}$ /s^2.
}
\end{ex}

% ===== Câu 91 =====
% ID: AIQ_L10C3_FULL_0901
% Mức: VD
\begin{ex}
% ID: AIQ_L10C3_FULL_0901
% Mức: VD
Vật $6\,\text{kg}$ trên sàn ngang có μ=0,3, chịu lực kéo ngang $24\,\text{N}$. Lấy g= $10\,\text{m}$ /s^2. Tính gia tốc.
\choice
{\True $1 m/s^2$}
{$2 m/s^2$}
{$0 m/s^2$}
{$2 m/s^2$}
\loigiai{
Chọn A. F_ms=μmg=18 $N; a=(F−F_ms)/m=(24−18)/6=1\,\text{m}$ /s^2.
}
\end{ex}

% ===== Câu 92 =====
% ID: AIQ_L10C3_FULL_0902
% Mức: TH
\begin{ex}
% ID: AIQ_L10C3_FULL_0902
% Mức: TH
Vật $5\,\text{kg}$ trên sàn ngang có μ=0,3, chịu lực kéo ngang $20\,\text{N}$. Lấy g= $10\,\text{m}$ /s^2. Tính gia tốc.
\choiceTF
{\True Giá trị cần tìm là $1\,\text{m}$ /s^2.}
{Giá trị cần tìm là $2\,\text{m}$ /s^2.}
{\True Phải xác định đúng các lực tác dụng và chọn chiều dương trước khi lập phương trình.}
{Có thể bỏ qua điều kiện áp dụng và chiều của lực mà kết quả vẫn luôn đúng.}
\loigiai{
\begin{itemchoice}
\item (Đúng) Giá trị cần tìm là $1\,\text{m}$ /s^2. (Vì): F_ms=μmg=15 $N; a=(F−F_ms)/m=(20−15)/5=1\,\text{m}$ /s^2. b) Sai vì sai giá trị. c) Đúng. d) Sai vì phải kiểm tra điều kiện áp dụng và chiều lực
\item (Sai) Giá trị cần tìm là $2\,\text{m}$ /s^2.
\item (Đúng) Phải xác định đúng các lực tác dụng và chọn chiều dương trước khi lập phương trình.
\item (Sai) Có thể bỏ qua điều kiện áp dụng và chiều của lực mà kết quả vẫn luôn đúng.
\end{itemchoice}
}
\end{ex}

% ===== Câu 93 =====
% ID: AIQ_L10C3_FULL_0903
% Mức: TH
\begin{ex}
% ID: AIQ_L10C3_FULL_0903
% Mức: TH
Vật $6\,\text{kg}$ trên sàn ngang có μ=0,4, chịu lực kéo ngang $36\,\text{N}$. Lấy g= $10\,\text{m}$ /s^2. Tính gia tốc.
\choiceTF
{\True Giá trị cần tìm là $2\,\text{m}$ /s^2.}
{Giá trị cần tìm là $3\,\text{m}$ /s^2.}
{\True Phải xác định đúng các lực tác dụng và chọn chiều dương trước khi lập phương trình.}
{Có thể bỏ qua điều kiện áp dụng và chiều của lực mà kết quả vẫn luôn đúng.}
\loigiai{
\begin{itemchoice}
\item (Đúng) Giá trị cần tìm là $2\,\text{m}$ /s^2. (Vì): F_ms=μmg=24 $N; a=(F−F_ms)/m=(36−24)/6=2\,\text{m}$ /s^2. b) Sai vì sai giá trị. c) Đúng. d) Sai vì phải kiểm tra điều kiện áp dụng và chiều lực
\item (Sai) Giá trị cần tìm là $3\,\text{m}$ /s^2.
\item (Đúng) Phải xác định đúng các lực tác dụng và chọn chiều dương trước khi lập phương trình.
\item (Sai) Có thể bỏ qua điều kiện áp dụng và chiều của lực mà kết quả vẫn luôn đúng.
\end{itemchoice}
}
\end{ex}

% ===== Câu 94 =====
% ID: AIQ_L10C3_FULL_0904
% Mức: VD
\begin{ex}
% ID: AIQ_L10C3_FULL_0904
% Mức: VD
Vật $2\,\text{kg}$ trên sàn ngang có μ=0,5, chịu lực kéo ngang $16\,\text{N}$. Lấy g= $10\,\text{m}$ /s^2. Tính gia tốc.
\choiceTF
{\True Giá trị cần tìm là $3\,\text{m}$ /s^2.}
{Giá trị cần tìm là $4\,\text{m}$ /s^2.}
{\True Phải xác định đúng các lực tác dụng và chọn chiều dương trước khi lập phương trình.}
{Có thể bỏ qua điều kiện áp dụng và chiều của lực mà kết quả vẫn luôn đúng.}
\loigiai{
\begin{itemchoice}
\item (Đúng) Giá trị cần tìm là $3\,\text{m}$ /s^2. (Vì): F_ms=μmg=10 $N; a=(F−F_ms)/m=(16−10)/2=3\,\text{m}$ /s^2. b) Sai vì sai giá trị. c) Đúng. d) Sai vì phải kiểm tra điều kiện áp dụng và chiều lực
\item (Sai) Giá trị cần tìm là $4\,\text{m}$ /s^2.
\item (Đúng) Phải xác định đúng các lực tác dụng và chọn chiều dương trước khi lập phương trình.
\item (Sai) Có thể bỏ qua điều kiện áp dụng và chiều của lực mà kết quả vẫn luôn đúng.
\end{itemchoice}
}
\end{ex}

% ===== Câu 95 =====
% ID: AIQ_L10C3_FULL_0905
% Mức: VD
\begin{ex}
% ID: AIQ_L10C3_FULL_0905
% Mức: VD
Vật $3\,\text{kg}$ trên sàn ngang có μ=0,2, chịu lực kéo ngang $18\,\text{N}$. Lấy g= $10\,\text{m}$ /s^2. Tính gia tốc.
\choiceTF
{\True Giá trị cần tìm là $4\,\text{m}$ /s^2.}
{Giá trị cần tìm là $5\,\text{m}$ /s^2.}
{\True Phải xác định đúng các lực tác dụng và chọn chiều dương trước khi lập phương trình.}
{Có thể bỏ qua điều kiện áp dụng và chiều của lực mà kết quả vẫn luôn đúng.}
\loigiai{
\begin{itemchoice}
\item (Đúng) Giá trị cần tìm là $4\,\text{m}$ /s^2. (Vì): F_ms=μmg=6 $N; a=(F−F_ms)/m=(18−6)/3=4\,\text{m}$ /s^2. b) Sai vì sai giá trị. c) Đúng. d) Sai vì phải kiểm tra điều kiện áp dụng và chiều lực
\item (Sai) Giá trị cần tìm là $5\,\text{m}$ /s^2.
\item (Đúng) Phải xác định đúng các lực tác dụng và chọn chiều dương trước khi lập phương trình.
\item (Sai) Có thể bỏ qua điều kiện áp dụng và chiều của lực mà kết quả vẫn luôn đúng.
\end{itemchoice}
}
\end{ex}

% ===== Câu 96 =====
% ID: AIQ_L10C3_FULL_0906
% Mức: VD
\begin{ex}
% ID: AIQ_L10C3_FULL_0906
% Mức: VD
Vật $4\,\text{kg}$ trên sàn ngang có μ=0,3, chịu lực kéo ngang $16\,\text{N}$. Lấy g= $10\,\text{m}$ /s^2. Tính gia tốc.
\choiceTF
{\True Giá trị cần tìm là $1\,\text{m}$ /s^2.}
{Giá trị cần tìm là $2\,\text{m}$ /s^2.}
{\True Phải xác định đúng các lực tác dụng và chọn chiều dương trước khi lập phương trình.}
{Có thể bỏ qua điều kiện áp dụng và chiều của lực mà kết quả vẫn luôn đúng.}
\loigiai{
\begin{itemchoice}
\item (Đúng) Giá trị cần tìm là $1\,\text{m}$ /s^2. (Vì): F_ms=μmg=12 $N; a=(F−F_ms)/m=(16−12)/4=1\,\text{m}$ /s^2. b) Sai vì sai giá trị. c) Đúng. d) Sai vì phải kiểm tra điều kiện áp dụng và chiều lực
\item (Sai) Giá trị cần tìm là $2\,\text{m}$ /s^2.
\item (Đúng) Phải xác định đúng các lực tác dụng và chọn chiều dương trước khi lập phương trình.
\item (Sai) Có thể bỏ qua điều kiện áp dụng và chiều của lực mà kết quả vẫn luôn đúng.
\end{itemchoice}
}
\end{ex}

% ===== Câu 97 =====
% ID: AIQ_L10C3_FULL_0907
% Mức: TH
\begin{ex}
% ID: AIQ_L10C3_FULL_0907
% Mức: TH
Vật $3\,\text{kg}$ trên sàn ngang có μ=0,4, chịu lực kéo ngang $18\,\text{N}$. Lấy g= $10\,\text{m}$ /s^2. Tính gia tốc.
\shortans{2}
\loigiai{
F_ms=μmg=12 $N; a=(F−F_ms)/m=(18−12)/3=2\,\text{m}$ /s^2.
}
\end{ex}

% ===== Câu 98 =====
% ID: AIQ_L10C3_FULL_0908
% Mức: TH
\begin{ex}
% ID: AIQ_L10C3_FULL_0908
% Mức: TH
Vật $4\,\text{kg}$ trên sàn ngang có μ=0,5, chịu lực kéo ngang $32\,\text{N}$. Lấy g= $10\,\text{m}$ /s^2. Tính gia tốc.
\shortans{3}
\loigiai{
F_ms=μmg=20 $N; a=(F−F_ms)/m=(32−20)/4=3\,\text{m}$ /s^2.
}
\end{ex}

% ===== Câu 99 =====
% ID: AIQ_L10C3_FULL_0909
% Mức: VD
\begin{ex}
% ID: AIQ_L10C3_FULL_0909
% Mức: VD
Vật $5\,\text{kg}$ trên sàn ngang có μ=0,2, chịu lực kéo ngang $30\,\text{N}$. Lấy g= $10\,\text{m}$ /s^2. Tính gia tốc.
\shortans{4}
\loigiai{
F_ms=μmg=10 $N; a=(F−F_ms)/m=(30−10)/5=4\,\text{m}$ /s^2.
}
\end{ex}

% ===== Câu 100 =====
% ID: AIQ_L10C3_FULL_0910
% Mức: VD
\begin{ex}
% ID: AIQ_L10C3_FULL_0910
% Mức: VD
Vật $6\,\text{kg}$ trên sàn ngang có μ=0,3, chịu lực kéo ngang $24\,\text{N}$. Lấy g= $10\,\text{m}$ /s^2. Tính gia tốc.
\shortans{1}
\loigiai{
F_ms=μmg=18 $N; a=(F−F_ms)/m=(24−18)/6=1\,\text{m}$ /s^2.
}
\end{ex}

% ===== Câu 101 =====
% ID: AIQ_L10C3_FULL_0911
% Mức: VD
\begin{ex}
% ID: AIQ_L10C3_FULL_0911
% Mức: VD
Vật $2\,\text{kg}$ trên sàn ngang có μ=0,4, chịu lực kéo ngang $12\,\text{N}$. Lấy g= $10\,\text{m}$ /s^2. Tính gia tốc.
\shortans{2}
\loigiai{
F_ms=μmg=8 $N; a=(F−F_ms)/m=(12−8)/2=2\,\text{m}$ /s^2.
}
\end{ex}

% ===== Câu 102 =====
% ID: AIQ_L10C3_FULL_0912
% Mức: VDC
\begin{ex}
% ID: AIQ_L10C3_FULL_0912
% Mức: VDC
Vật $3\,\text{kg}$ trên sàn ngang có μ=0,5, chịu lực kéo ngang $24\,\text{N}$. Lấy g= $10\,\text{m}$ /s^2. Tính gia tốc.
\shortans{3}
\loigiai{
F_ms=μmg=15 $N; a=(F−F_ms)/m=(24−15)/3=3\,\text{m}$ /s^2.
}
\end{ex}

% ===== Câu 103 =====
% ID: AIQ_L10C3_FULL_0913
% Mức: VD
\begin{bt}
% ID: AIQ_L10C3_FULL_0913
% Mức: VD
Vật $6\,\text{kg}$ trên sàn ngang có μ=0,5, chịu lực kéo ngang $48\,\text{N}$. Lấy g= $10\,\text{m}$ /s^2. Tính gia tốc.
\loigiai{
Phân tích lực, chọn hệ quy chiếu và áp dụng định luật hoặc quy tắc phù hợp. F_ms=μmg=30 $N; a=(F−F_ms)/m=(48−30)/6=3\,\text{m}$ /s^2.
}
\end{bt}

% ===== Câu 104 =====
% ID: AIQ_L10C3_FULL_0914
% Mức: VD
\begin{bt}
% ID: AIQ_L10C3_FULL_0914
% Mức: VD
Vật $2\,\text{kg}$ trên sàn ngang có μ=0,2, chịu lực kéo ngang $12\,\text{N}$. Lấy g= $10\,\text{m}$ /s^2. Tính gia tốc.
\loigiai{
Phân tích lực, chọn hệ quy chiếu và áp dụng định luật hoặc quy tắc phù hợp. F_ms=μmg=4 $N; a=(F−F_ms)/m=(12−4)/2=4\,\text{m}$ /s^2.
}
\end{bt}

% ===== Câu 105 =====
% ID: AIQ_L10C3_FULL_0915
% Mức: VDC
\begin{bt}
% ID: AIQ_L10C3_FULL_0915
% Mức: VDC
Vật $3\,\text{kg}$ trên sàn ngang có μ=0,3, chịu lực kéo ngang $12\,\text{N}$. Lấy g= $10\,\text{m}$ /s^2. Tính gia tốc.
\loigiai{
Phân tích lực, chọn hệ quy chiếu và áp dụng định luật hoặc quy tắc phù hợp. F_ms=μmg=9 $N; a=(F−F_ms)/m=(12−9)/3=1\,\text{m}$ /s^2.
}
\end{bt}

% ===== Câu 106 =====
% ID: AIQ_L10C3_FULL_0916
% Mức: VDC
\begin{bt}
% ID: AIQ_L10C3_FULL_0916
% Mức: VDC
Vật $4\,\text{kg}$ trên sàn ngang có μ=0,4, chịu lực kéo ngang $24\,\text{N}$. Lấy g= $10\,\text{m}$ /s^2. Tính gia tốc.
\loigiai{
Phân tích lực, chọn hệ quy chiếu và áp dụng định luật hoặc quy tắc phù hợp. F_ms=μmg=16 $N; a=(F−F_ms)/m=(24−16)/4=2\,\text{m}$ /s^2.
}
\end{bt}

% ===== Câu 107 =====
% ID: AIQ_L10C3_FULL_0917
% Mức: VDC
\begin{bt}
% ID: AIQ_L10C3_FULL_0917
% Mức: VDC
Vật $5\,\text{kg}$ trên sàn ngang có μ=0,5, chịu lực kéo ngang $40\,\text{N}$. Lấy g= $10\,\text{m}$ /s^2. Tính gia tốc.
\loigiai{
Phân tích lực, chọn hệ quy chiếu và áp dụng định luật hoặc quy tắc phù hợp. F_ms=μmg=25 $N; a=(F−F_ms)/m=(40−25)/5=3\,\text{m}$ /s^2.
}
\end{bt}

% ===== Câu 108 =====
% ID: AIQ_L10C3_FULL_0918
% Mức: VDC
\begin{bt}
% ID: AIQ_L10C3_FULL_0918
% Mức: VDC
Vật $6\,\text{kg}$ trên sàn ngang có μ=0,2, chịu lực kéo ngang $36\,\text{N}$. Lấy g= $10\,\text{m}$ /s^2. Tính gia tốc.
\loigiai{
Phân tích lực, chọn hệ quy chiếu và áp dụng định luật hoặc quy tắc phù hợp. F_ms=μmg=12 $N; a=(F−F_ms)/m=(36−12)/6=4\,\text{m}$ /s^2.
}
\end{bt}

% ===== Câu 109 =====
% ID: AIQ_L10C3_FULL_0919
% Mức: NB
\dangbt{Chuyển động trên mặt phẳng nghiêng có ma sát}
\begin{ex}
% ID: AIQ_L10C3_FULL_0919
% Mức: NB
Vật $2\,\text{kg}$ trượt xuống mặt phẳng nghiêng 37°, có μ=0,25; lấy g= $10\,\text{m}$ /s^2, sin37°=0,6, cos37°=0,8. Tính gia tốc.
\choice
{\True $4 m/s^2$}
{$5 m/s^2$}
{$3 m/s^2$}
{$8 m/s^2$}
\loigiai{
Chọn A. a=g $(sinα−μcosα)=10(0,6−0,25·0,8)=4\,\text{m}$ /s^2.
}
\end{ex}

% ===== Câu 110 =====
% ID: AIQ_L10C3_FULL_0920
% Mức: NB
\begin{ex}
% ID: AIQ_L10C3_FULL_0920
% Mức: NB
Vật $3\,\text{kg}$ trượt xuống mặt phẳng nghiêng 37°, có μ=0,25; lấy g= $10\,\text{m}$ /s^2, sin37°=0,6, cos37°=0,8. Tính gia tốc.
\choice
{$4 m/s^2$}
{\True $4 m/s^2$}
{$3 m/s^2$}
{$8 m/s^2$}
\loigiai{
Chọn B. a=g $(sinα−μcosα)=10(0,6−0,25·0,8)=4\,\text{m}$ /s^2.
}
\end{ex}

% ===== Câu 111 =====
% ID: AIQ_L10C3_FULL_0921
% Mức: NB
\begin{ex}
% ID: AIQ_L10C3_FULL_0921
% Mức: NB
Vật $4\,\text{kg}$ trượt xuống mặt phẳng nghiêng 37°, có μ=0,25; lấy g= $10\,\text{m}$ /s^2, sin37°=0,6, cos37°=0,8. Tính gia tốc.
\choice
{$4 m/s^2$}
{$5 m/s^2$}
{\True $4 m/s^2$}
{$8 m/s^2$}
\loigiai{
Chọn C. a=g $(sinα−μcosα)=10(0,6−0,25·0,8)=4\,\text{m}$ /s^2.
}
\end{ex}

% ===== Câu 112 =====
% ID: AIQ_L10C3_FULL_0922
% Mức: NB
\begin{ex}
% ID: AIQ_L10C3_FULL_0922
% Mức: NB
Vật $5\,\text{kg}$ trượt xuống mặt phẳng nghiêng 37°, có μ=0,25; lấy g= $10\,\text{m}$ /s^2, sin37°=0,6, cos37°=0,8. Tính gia tốc.
\choice
{$4 m/s^2$}
{$5 m/s^2$}
{$3 m/s^2$}
{\True $4 m/s^2$}
\loigiai{
Chọn D. a=g $(sinα−μcosα)=10(0,6−0,25·0,8)=4\,\text{m}$ /s^2.
}
\end{ex}

% ===== Câu 113 =====
% ID: AIQ_L10C3_FULL_0923
% Mức: TH
\begin{ex}
% ID: AIQ_L10C3_FULL_0923
% Mức: TH
Vật $6\,\text{kg}$ trượt xuống mặt phẳng nghiêng 37°, có μ=0,25; lấy g= $10\,\text{m}$ /s^2, sin37°=0,6, cos37°=0,8. Tính gia tốc.
\choice
{\True $4 m/s^2$}
{$5 m/s^2$}
{$3 m/s^2$}
{$8 m/s^2$}
\loigiai{
Chọn A. a=g $(sinα−μcosα)=10(0,6−0,25·0,8)=4\,\text{m}$ /s^2.
}
\end{ex}

% ===== Câu 114 =====
% ID: AIQ_L10C3_FULL_0924
% Mức: TH
\begin{ex}
% ID: AIQ_L10C3_FULL_0924
% Mức: TH
Vật $2\,\text{kg}$ trượt xuống mặt phẳng nghiêng 37°, có μ=0,25; lấy g= $10\,\text{m}$ /s^2, sin37°=0,6, cos37°=0,8. Tính gia tốc.
\choice
{$4 m/s^2$}
{\True $4 m/s^2$}
{$3 m/s^2$}
{$8 m/s^2$}
\loigiai{
Chọn B. a=g $(sinα−μcosα)=10(0,6−0,25·0,8)=4\,\text{m}$ /s^2.
}
\end{ex}

% ===== Câu 115 =====
% ID: AIQ_L10C3_FULL_0925
% Mức: TH
\begin{ex}
% ID: AIQ_L10C3_FULL_0925
% Mức: TH
Vật $3\,\text{kg}$ trượt xuống mặt phẳng nghiêng 37°, có μ=0,25; lấy g= $10\,\text{m}$ /s^2, sin37°=0,6, cos37°=0,8. Tính gia tốc.
\choice
{$4 m/s^2$}
{$5 m/s^2$}
{\True $4 m/s^2$}
{$8 m/s^2$}
\loigiai{
Chọn C. a=g $(sinα−μcosα)=10(0,6−0,25·0,8)=4\,\text{m}$ /s^2.
}
\end{ex}

% ===== Câu 116 =====
% ID: AIQ_L10C3_FULL_0926
% Mức: TH
\begin{ex}
% ID: AIQ_L10C3_FULL_0926
% Mức: TH
Vật $4\,\text{kg}$ trượt xuống mặt phẳng nghiêng 37°, có μ=0,25; lấy g= $10\,\text{m}$ /s^2, sin37°=0,6, cos37°=0,8. Tính gia tốc.
\choice
{$4 m/s^2$}
{$5 m/s^2$}
{$3 m/s^2$}
{\True $4 m/s^2$}
\loigiai{
Chọn D. a=g $(sinα−μcosα)=10(0,6−0,25·0,8)=4\,\text{m}$ /s^2.
}
\end{ex}

% ===== Câu 117 =====
% ID: AIQ_L10C3_FULL_0927
% Mức: VD
\begin{ex}
% ID: AIQ_L10C3_FULL_0927
% Mức: VD
Vật $5\,\text{kg}$ trượt xuống mặt phẳng nghiêng 37°, có μ=0,25; lấy g= $10\,\text{m}$ /s^2, sin37°=0,6, cos37°=0,8. Tính gia tốc.
\choice
{\True $4 m/s^2$}
{$5 m/s^2$}
{$3 m/s^2$}
{$8 m/s^2$}
\loigiai{
Chọn A. a=g $(sinα−μcosα)=10(0,6−0,25·0,8)=4\,\text{m}$ /s^2.
}
\end{ex}

% ===== Câu 118 =====
% ID: AIQ_L10C3_FULL_0928
% Mức: VD
\begin{ex}
% ID: AIQ_L10C3_FULL_0928
% Mức: VD
Vật $6\,\text{kg}$ trượt xuống mặt phẳng nghiêng 37°, có μ=0,25; lấy g= $10\,\text{m}$ /s^2, sin37°=0,6, cos37°=0,8. Tính gia tốc.
\choice
{$4 m/s^2$}
{\True $4 m/s^2$}
{$3 m/s^2$}
{$8 m/s^2$}
\loigiai{
Chọn B. a=g $(sinα−μcosα)=10(0,6−0,25·0,8)=4\,\text{m}$ /s^2.
}
\end{ex}

% ===== Câu 119 =====
% ID: AIQ_L10C3_FULL_0929
% Mức: TH
\begin{ex}
% ID: AIQ_L10C3_FULL_0929
% Mức: TH
Vật $5\,\text{kg}$ trượt xuống mặt phẳng nghiêng 37°, có μ=0,25; lấy g= $10\,\text{m}$ /s^2, sin37°=0,6, cos37°=0,8. Tính gia tốc.
\choiceTF
{\True Giá trị cần tìm là $4\,\text{m}$ /s^2.}
{Giá trị cần tìm là $5\,\text{m}$ /s^2.}
{\True Phải xác định đúng các lực tác dụng và chọn chiều dương trước khi lập phương trình.}
{Có thể bỏ qua điều kiện áp dụng và chiều của lực mà kết quả vẫn luôn đúng.}
\loigiai{
\begin{itemchoice}
\item (Đúng) Giá trị cần tìm là $4\,\text{m}$ /s^2. (Vì): a=g $(sinα−μcosα)=10(0,6−0,25·0,8)=4\,\text{m}$ /s^2. b) Sai vì sai giá trị. c) Đúng. d) Sai vì phải kiểm tra điều kiện áp dụng và chiều lực
\item (Sai) Giá trị cần tìm là $5\,\text{m}$ /s^2.
\item (Đúng) Phải xác định đúng các lực tác dụng và chọn chiều dương trước khi lập phương trình.
\item (Sai) Có thể bỏ qua điều kiện áp dụng và chiều của lực mà kết quả vẫn luôn đúng.
\end{itemchoice}
}
\end{ex}

% ===== Câu 120 =====
% ID: AIQ_L10C3_FULL_0930
% Mức: TH
\begin{ex}
% ID: AIQ_L10C3_FULL_0930
% Mức: TH
Vật $6\,\text{kg}$ trượt xuống mặt phẳng nghiêng 37°, có μ=0,25; lấy g= $10\,\text{m}$ /s^2, sin37°=0,6, cos37°=0,8. Tính gia tốc.
\choiceTF
{\True Giá trị cần tìm là $4\,\text{m}$ /s^2.}
{Giá trị cần tìm là $5\,\text{m}$ /s^2.}
{\True Phải xác định đúng các lực tác dụng và chọn chiều dương trước khi lập phương trình.}
{Có thể bỏ qua điều kiện áp dụng và chiều của lực mà kết quả vẫn luôn đúng.}
\loigiai{
\begin{itemchoice}
\item (Đúng) Giá trị cần tìm là $4\,\text{m}$ /s^2. (Vì): a=g $(sinα−μcosα)=10(0,6−0,25·0,8)=4\,\text{m}$ /s^2. b) Sai vì sai giá trị. c) Đúng. d) Sai vì phải kiểm tra điều kiện áp dụng và chiều lực
\item (Sai) Giá trị cần tìm là $5\,\text{m}$ /s^2.
\item (Đúng) Phải xác định đúng các lực tác dụng và chọn chiều dương trước khi lập phương trình.
\item (Sai) Có thể bỏ qua điều kiện áp dụng và chiều của lực mà kết quả vẫn luôn đúng.
\end{itemchoice}
}
\end{ex}

% ===== Câu 121 =====
% ID: AIQ_L10C3_FULL_0931
% Mức: VD
\begin{ex}
% ID: AIQ_L10C3_FULL_0931
% Mức: VD
Vật $2\,\text{kg}$ trượt xuống mặt phẳng nghiêng 37°, có μ=0,25; lấy g= $10\,\text{m}$ /s^2, sin37°=0,6, cos37°=0,8. Tính gia tốc.
\choiceTF
{\True Giá trị cần tìm là $4\,\text{m}$ /s^2.}
{Giá trị cần tìm là $5\,\text{m}$ /s^2.}
{\True Phải xác định đúng các lực tác dụng và chọn chiều dương trước khi lập phương trình.}
{Có thể bỏ qua điều kiện áp dụng và chiều của lực mà kết quả vẫn luôn đúng.}
\loigiai{
\begin{itemchoice}
\item (Đúng) Giá trị cần tìm là $4\,\text{m}$ /s^2. (Vì): a=g $(sinα−μcosα)=10(0,6−0,25·0,8)=4\,\text{m}$ /s^2. b) Sai vì sai giá trị. c) Đúng. d) Sai vì phải kiểm tra điều kiện áp dụng và chiều lực
\item (Sai) Giá trị cần tìm là $5\,\text{m}$ /s^2.
\item (Đúng) Phải xác định đúng các lực tác dụng và chọn chiều dương trước khi lập phương trình.
\item (Sai) Có thể bỏ qua điều kiện áp dụng và chiều của lực mà kết quả vẫn luôn đúng.
\end{itemchoice}
}
\end{ex}

% ===== Câu 122 =====
% ID: AIQ_L10C3_FULL_0932
% Mức: VD
\begin{ex}
% ID: AIQ_L10C3_FULL_0932
% Mức: VD
Vật $3\,\text{kg}$ trượt xuống mặt phẳng nghiêng 37°, có μ=0,25; lấy g= $10\,\text{m}$ /s^2, sin37°=0,6, cos37°=0,8. Tính gia tốc.
\choiceTF
{\True Giá trị cần tìm là $4\,\text{m}$ /s^2.}
{Giá trị cần tìm là $5\,\text{m}$ /s^2.}
{\True Phải xác định đúng các lực tác dụng và chọn chiều dương trước khi lập phương trình.}
{Có thể bỏ qua điều kiện áp dụng và chiều của lực mà kết quả vẫn luôn đúng.}
\loigiai{
\begin{itemchoice}
\item (Đúng) Giá trị cần tìm là $4\,\text{m}$ /s^2. (Vì): a=g $(sinα−μcosα)=10(0,6−0,25·0,8)=4\,\text{m}$ /s^2. b) Sai vì sai giá trị. c) Đúng. d) Sai vì phải kiểm tra điều kiện áp dụng và chiều lực
\item (Sai) Giá trị cần tìm là $5\,\text{m}$ /s^2.
\item (Đúng) Phải xác định đúng các lực tác dụng và chọn chiều dương trước khi lập phương trình.
\item (Sai) Có thể bỏ qua điều kiện áp dụng và chiều của lực mà kết quả vẫn luôn đúng.
\end{itemchoice}
}
\end{ex}

% ===== Câu 123 =====
% ID: AIQ_L10C3_FULL_0933
% Mức: VD
\begin{ex}
% ID: AIQ_L10C3_FULL_0933
% Mức: VD
Vật $4\,\text{kg}$ trượt xuống mặt phẳng nghiêng 37°, có μ=0,25; lấy g= $10\,\text{m}$ /s^2, sin37°=0,6, cos37°=0,8. Tính gia tốc.
\choiceTF
{\True Giá trị cần tìm là $4\,\text{m}$ /s^2.}
{Giá trị cần tìm là $5\,\text{m}$ /s^2.}
{\True Phải xác định đúng các lực tác dụng và chọn chiều dương trước khi lập phương trình.}
{Có thể bỏ qua điều kiện áp dụng và chiều của lực mà kết quả vẫn luôn đúng.}
\loigiai{
\begin{itemchoice}
\item (Đúng) Giá trị cần tìm là $4\,\text{m}$ /s^2. (Vì): a=g $(sinα−μcosα)=10(0,6−0,25·0,8)=4\,\text{m}$ /s^2. b) Sai vì sai giá trị. c) Đúng. d) Sai vì phải kiểm tra điều kiện áp dụng và chiều lực
\item (Sai) Giá trị cần tìm là $5\,\text{m}$ /s^2.
\item (Đúng) Phải xác định đúng các lực tác dụng và chọn chiều dương trước khi lập phương trình.
\item (Sai) Có thể bỏ qua điều kiện áp dụng và chiều của lực mà kết quả vẫn luôn đúng.
\end{itemchoice}
}
\end{ex}

% ===== Câu 124 =====
% ID: AIQ_L10C3_FULL_0934
% Mức: TH
\begin{ex}
% ID: AIQ_L10C3_FULL_0934
% Mức: TH
Vật $3\,\text{kg}$ trượt xuống mặt phẳng nghiêng 37°, có μ=0,25; lấy g= $10\,\text{m}$ /s^2, sin37°=0,6, cos37°=0,8. Tính gia tốc.
\shortans{4}
\loigiai{
a=g $(sinα−μcosα)=10(0,6−0,25·0,8)=4\,\text{m}$ /s^2.
}
\end{ex}

% ===== Câu 125 =====
% ID: AIQ_L10C3_FULL_0935
% Mức: TH
\begin{ex}
% ID: AIQ_L10C3_FULL_0935
% Mức: TH
Vật $4\,\text{kg}$ trượt xuống mặt phẳng nghiêng 37°, có μ=0,25; lấy g= $10\,\text{m}$ /s^2, sin37°=0,6, cos37°=0,8. Tính gia tốc.
\shortans{4}
\loigiai{
a=g $(sinα−μcosα)=10(0,6−0,25·0,8)=4\,\text{m}$ /s^2.
}
\end{ex}

% ===== Câu 126 =====
% ID: AIQ_L10C3_FULL_0936
% Mức: VD
\begin{ex}
% ID: AIQ_L10C3_FULL_0936
% Mức: VD
Vật $5\,\text{kg}$ trượt xuống mặt phẳng nghiêng 37°, có μ=0,25; lấy g= $10\,\text{m}$ /s^2, sin37°=0,6, cos37°=0,8. Tính gia tốc.
\shortans{4}
\loigiai{
a=g $(sinα−μcosα)=10(0,6−0,25·0,8)=4\,\text{m}$ /s^2.
}
\end{ex}

% ===== Câu 127 =====
% ID: AIQ_L10C3_FULL_0937
% Mức: VD
\begin{ex}
% ID: AIQ_L10C3_FULL_0937
% Mức: VD
Vật $6\,\text{kg}$ trượt xuống mặt phẳng nghiêng 37°, có μ=0,25; lấy g= $10\,\text{m}$ /s^2, sin37°=0,6, cos37°=0,8. Tính gia tốc.
\shortans{4}
\loigiai{
a=g $(sinα−μcosα)=10(0,6−0,25·0,8)=4\,\text{m}$ /s^2.
}
\end{ex}

% ===== Câu 128 =====
% ID: AIQ_L10C3_FULL_0938
% Mức: VD
\begin{ex}
% ID: AIQ_L10C3_FULL_0938
% Mức: VD
Vật $2\,\text{kg}$ trượt xuống mặt phẳng nghiêng 37°, có μ=0,25; lấy g= $10\,\text{m}$ /s^2, sin37°=0,6, cos37°=0,8. Tính gia tốc.
\shortans{4}
\loigiai{
a=g $(sinα−μcosα)=10(0,6−0,25·0,8)=4\,\text{m}$ /s^2.
}
\end{ex}

% ===== Câu 129 =====
% ID: AIQ_L10C3_FULL_0939
% Mức: VDC
\begin{ex}
% ID: AIQ_L10C3_FULL_0939
% Mức: VDC
Vật $3\,\text{kg}$ trượt xuống mặt phẳng nghiêng 37°, có μ=0,25; lấy g= $10\,\text{m}$ /s^2, sin37°=0,6, cos37°=0,8. Tính gia tốc.
\shortans{4}
\loigiai{
a=g $(sinα−μcosα)=10(0,6−0,25·0,8)=4\,\text{m}$ /s^2.
}
\end{ex}

% ===== Câu 130 =====
% ID: AIQ_L10C3_FULL_0940
% Mức: VD
\begin{bt}
% ID: AIQ_L10C3_FULL_0940
% Mức: VD
Vật $6\,\text{kg}$ trượt xuống mặt phẳng nghiêng 37°, có μ=0,25; lấy g= $10\,\text{m}$ /s^2, sin37°=0,6, cos37°=0,8. Tính gia tốc.
\loigiai{
Phân tích lực, chọn hệ quy chiếu và áp dụng định luật hoặc quy tắc phù hợp. a=g $(sinα−μcosα)=10(0,6−0,25·0,8)=4\,\text{m}$ /s^2.
}
\end{bt}

% ===== Câu 131 =====
% ID: AIQ_L10C3_FULL_0941
% Mức: VD
\begin{bt}
% ID: AIQ_L10C3_FULL_0941
% Mức: VD
Vật $2\,\text{kg}$ trượt xuống mặt phẳng nghiêng 37°, có μ=0,25; lấy g= $10\,\text{m}$ /s^2, sin37°=0,6, cos37°=0,8. Tính gia tốc.
\loigiai{
Phân tích lực, chọn hệ quy chiếu và áp dụng định luật hoặc quy tắc phù hợp. a=g $(sinα−μcosα)=10(0,6−0,25·0,8)=4\,\text{m}$ /s^2.
}
\end{bt}

% ===== Câu 132 =====
% ID: AIQ_L10C3_FULL_0942
% Mức: VDC
\begin{bt}
% ID: AIQ_L10C3_FULL_0942
% Mức: VDC
Vật $3\,\text{kg}$ trượt xuống mặt phẳng nghiêng 37°, có μ=0,25; lấy g= $10\,\text{m}$ /s^2, sin37°=0,6, cos37°=0,8. Tính gia tốc.
\loigiai{
Phân tích lực, chọn hệ quy chiếu và áp dụng định luật hoặc quy tắc phù hợp. a=g $(sinα−μcosα)=10(0,6−0,25·0,8)=4\,\text{m}$ /s^2.
}
\end{bt}

% ===== Câu 133 =====
% ID: AIQ_L10C3_FULL_0943
% Mức: VDC
\begin{bt}
% ID: AIQ_L10C3_FULL_0943
% Mức: VDC
Vật $4\,\text{kg}$ trượt xuống mặt phẳng nghiêng 37°, có μ=0,25; lấy g= $10\,\text{m}$ /s^2, sin37°=0,6, cos37°=0,8. Tính gia tốc.
\loigiai{
Phân tích lực, chọn hệ quy chiếu và áp dụng định luật hoặc quy tắc phù hợp. a=g $(sinα−μcosα)=10(0,6−0,25·0,8)=4\,\text{m}$ /s^2.
}
\end{bt}

% ===== Câu 134 =====
% ID: AIQ_L10C3_FULL_0944
% Mức: VDC
\begin{bt}
% ID: AIQ_L10C3_FULL_0944
% Mức: VDC
Vật $5\,\text{kg}$ trượt xuống mặt phẳng nghiêng 37°, có μ=0,25; lấy g= $10\,\text{m}$ /s^2, sin37°=0,6, cos37°=0,8. Tính gia tốc.
\loigiai{
Phân tích lực, chọn hệ quy chiếu và áp dụng định luật hoặc quy tắc phù hợp. a=g $(sinα−μcosα)=10(0,6−0,25·0,8)=4\,\text{m}$ /s^2.
}
\end{bt}

% ===== Câu 135 =====
% ID: AIQ_L10C3_FULL_0945
% Mức: VDC
\begin{bt}
% ID: AIQ_L10C3_FULL_0945
% Mức: VDC
Vật $6\,\text{kg}$ trượt xuống mặt phẳng nghiêng 37°, có μ=0,25; lấy g= $10\,\text{m}$ /s^2, sin37°=0,6, cos37°=0,8. Tính gia tốc.
\loigiai{
Phân tích lực, chọn hệ quy chiếu và áp dụng định luật hoặc quy tắc phù hợp. a=g $(sinα−μcosα)=10(0,6−0,25·0,8)=4\,\text{m}$ /s^2.
}
\end{bt}

% ===== Câu 136 =====
% ID: AIQ_L10C3_FULL_0946
% Mức: NB
\dangbt{Bài toán tổng hợp và ứng dụng thực tiễn của lực ma sát}
\begin{ex}
% ID: AIQ_L10C3_FULL_0946
% Mức: NB
Muốn thùng $2\,\text{kg}$ trên sàn ngang chuyển động đều, hệ số ma sát trượt 0,2. Lấy g= $10\,\text{m}$ /s^2. Lực kéo ngang tối thiểu khi đã trượt là bao nhiêu?
\choice
{4 $N$}
{\True 4 $N$}
{3 $N$}
{8 $N$}
\loigiai{
Chọn B. Chuyển động đều nên lực kéo cân bằng ma sát: F=μmg= $4\,\text{N}$.
}
\end{ex}

% ===== Câu 137 =====
% ID: AIQ_L10C3_FULL_0947
% Mức: NB
\begin{ex}
% ID: AIQ_L10C3_FULL_0947
% Mức: NB
Muốn thùng $3\,\text{kg}$ trên sàn ngang chuyển động đều, hệ số ma sát trượt 0,3. Lấy g= $10\,\text{m}$ /s^2. Lực kéo ngang tối thiểu khi đã trượt là bao nhiêu?
\choice
{9 $N$}
{10 $N$}
{\True 9 $N$}
{18 $N$}
\loigiai{
Chọn C. Chuyển động đều nên lực kéo cân bằng ma sát: F=μmg= $9\,\text{N}$.
}
\end{ex}

% ===== Câu 138 =====
% ID: AIQ_L10C3_FULL_0948
% Mức: NB
\begin{ex}
% ID: AIQ_L10C3_FULL_0948
% Mức: NB
Muốn thùng $4\,\text{kg}$ trên sàn ngang chuyển động đều, hệ số ma sát trượt 0,4. Lấy g= $10\,\text{m}$ /s^2. Lực kéo ngang tối thiểu khi đã trượt là bao nhiêu?
\choice
{16 $N$}
{17 $N$}
{15 $N$}
{\True 16 $N$}
\loigiai{
Chọn D. Chuyển động đều nên lực kéo cân bằng ma sát: F=μmg= $16\,\text{N}$.
}
\end{ex}

% ===== Câu 139 =====
% ID: AIQ_L10C3_FULL_0949
% Mức: NB
\begin{ex}
% ID: AIQ_L10C3_FULL_0949
% Mức: NB
Muốn thùng $5\,\text{kg}$ trên sàn ngang chuyển động đều, hệ số ma sát trượt 0,5. Lấy g= $10\,\text{m}$ /s^2. Lực kéo ngang tối thiểu khi đã trượt là bao nhiêu?
\choice
{\True 25 $N$}
{26 $N$}
{24 $N$}
{50 $N$}
\loigiai{
Chọn A. Chuyển động đều nên lực kéo cân bằng ma sát: F=μmg= $25\,\text{N}$.
}
\end{ex}

% ===== Câu 140 =====
% ID: AIQ_L10C3_FULL_0950
% Mức: TH
\begin{ex}
% ID: AIQ_L10C3_FULL_0950
% Mức: TH
Muốn thùng $6\,\text{kg}$ trên sàn ngang chuyển động đều, hệ số ma sát trượt 0,2. Lấy g= $10\,\text{m}$ /s^2. Lực kéo ngang tối thiểu khi đã trượt là bao nhiêu?
\choice
{12 $N$}
{\True 12 $N$}
{11 $N$}
{24 $N$}
\loigiai{
Chọn B. Chuyển động đều nên lực kéo cân bằng ma sát: F=μmg= $12\,\text{N}$.
}
\end{ex}

% ===== Câu 141 =====
% ID: AIQ_L10C3_FULL_0951
% Mức: TH
\begin{ex}
% ID: AIQ_L10C3_FULL_0951
% Mức: TH
Muốn thùng $2\,\text{kg}$ trên sàn ngang chuyển động đều, hệ số ma sát trượt 0,3. Lấy g= $10\,\text{m}$ /s^2. Lực kéo ngang tối thiểu khi đã trượt là bao nhiêu?
\choice
{6 $N$}
{7 $N$}
{\True 6 $N$}
{12 $N$}
\loigiai{
Chọn C. Chuyển động đều nên lực kéo cân bằng ma sát: F=μmg= $6\,\text{N}$.
}
\end{ex}

% ===== Câu 142 =====
% ID: AIQ_L10C3_FULL_0952
% Mức: TH
\begin{ex}
% ID: AIQ_L10C3_FULL_0952
% Mức: TH
Muốn thùng $3\,\text{kg}$ trên sàn ngang chuyển động đều, hệ số ma sát trượt 0,4. Lấy g= $10\,\text{m}$ /s^2. Lực kéo ngang tối thiểu khi đã trượt là bao nhiêu?
\choice
{12 $N$}
{13 $N$}
{11 $N$}
{\True 12 $N$}
\loigiai{
Chọn D. Chuyển động đều nên lực kéo cân bằng ma sát: F=μmg= $12\,\text{N}$.
}
\end{ex}

% ===== Câu 143 =====
% ID: AIQ_L10C3_FULL_0953
% Mức: TH
\begin{ex}
% ID: AIQ_L10C3_FULL_0953
% Mức: TH
Muốn thùng $4\,\text{kg}$ trên sàn ngang chuyển động đều, hệ số ma sát trượt 0,5. Lấy g= $10\,\text{m}$ /s^2. Lực kéo ngang tối thiểu khi đã trượt là bao nhiêu?
\choice
{\True 20 $N$}
{21 $N$}
{19 $N$}
{40 $N$}
\loigiai{
Chọn A. Chuyển động đều nên lực kéo cân bằng ma sát: F=μmg= $20\,\text{N}$.
}
\end{ex}

% ===== Câu 144 =====
% ID: AIQ_L10C3_FULL_0954
% Mức: VD
\begin{ex}
% ID: AIQ_L10C3_FULL_0954
% Mức: VD
Muốn thùng $5\,\text{kg}$ trên sàn ngang chuyển động đều, hệ số ma sát trượt 0,2. Lấy g= $10\,\text{m}$ /s^2. Lực kéo ngang tối thiểu khi đã trượt là bao nhiêu?
\choice
{10 $N$}
{\True 10 $N$}
{9 $N$}
{20 $N$}
\loigiai{
Chọn B. Chuyển động đều nên lực kéo cân bằng ma sát: F=μmg= $10\,\text{N}$.
}
\end{ex}

% ===== Câu 145 =====
% ID: AIQ_L10C3_FULL_0955
% Mức: VD
\begin{ex}
% ID: AIQ_L10C3_FULL_0955
% Mức: VD
Muốn thùng $6\,\text{kg}$ trên sàn ngang chuyển động đều, hệ số ma sát trượt 0,3. Lấy g= $10\,\text{m}$ /s^2. Lực kéo ngang tối thiểu khi đã trượt là bao nhiêu?
\choice
{18 $N$}
{19 $N$}
{\True 18 $N$}
{36 $N$}
\loigiai{
Chọn C. Chuyển động đều nên lực kéo cân bằng ma sát: F=μmg= $18\,\text{N}$.
}
\end{ex}

% ===== Câu 146 =====
% ID: AIQ_L10C3_FULL_0956
% Mức: TH
\begin{ex}
% ID: AIQ_L10C3_FULL_0956
% Mức: TH
Muốn thùng $5\,\text{kg}$ trên sàn ngang chuyển động đều, hệ số ma sát trượt 0,3. Lấy g= $10\,\text{m}$ /s^2. Lực kéo ngang tối thiểu khi đã trượt là bao nhiêu?
\choiceTF
{\True Giá trị cần tìm là $15\,\text{N}$.}
{Giá trị cần tìm là $16\,\text{N}$.}
{\True Phải xác định đúng các lực tác dụng và chọn chiều dương trước khi lập phương trình.}
{Có thể bỏ qua điều kiện áp dụng và chiều của lực mà kết quả vẫn luôn đúng.}
\loigiai{
\begin{itemchoice}
\item (Đúng) Giá trị cần tìm là $15\,\text{N}$. (Vì): Chuyển động đều nên lực kéo cân bằng ma sát: F=μmg= $15\,\text{N}$. b) Sai vì sai giá trị. c) Đúng. d) Sai vì phải kiểm tra điều kiện áp dụng và chiều lực
\item (Sai) Giá trị cần tìm là $16\,\text{N}$.
\item (Đúng) Phải xác định đúng các lực tác dụng và chọn chiều dương trước khi lập phương trình.
\item (Sai) Có thể bỏ qua điều kiện áp dụng và chiều của lực mà kết quả vẫn luôn đúng.
\end{itemchoice}
}
\end{ex}

% ===== Câu 147 =====
% ID: AIQ_L10C3_FULL_0957
% Mức: TH
\begin{ex}
% ID: AIQ_L10C3_FULL_0957
% Mức: TH
Muốn thùng $6\,\text{kg}$ trên sàn ngang chuyển động đều, hệ số ma sát trượt 0,4. Lấy g= $10\,\text{m}$ /s^2. Lực kéo ngang tối thiểu khi đã trượt là bao nhiêu?
\choiceTF
{\True Giá trị cần tìm là $24\,\text{N}$.}
{Giá trị cần tìm là $25\,\text{N}$.}
{\True Phải xác định đúng các lực tác dụng và chọn chiều dương trước khi lập phương trình.}
{Có thể bỏ qua điều kiện áp dụng và chiều của lực mà kết quả vẫn luôn đúng.}
\loigiai{
\begin{itemchoice}
\item (Đúng) Giá trị cần tìm là $24\,\text{N}$. (Vì): Chuyển động đều nên lực kéo cân bằng ma sát: F=μmg= $24\,\text{N}$. b) Sai vì sai giá trị. c) Đúng. d) Sai vì phải kiểm tra điều kiện áp dụng và chiều lực
\item (Sai) Giá trị cần tìm là $25\,\text{N}$.
\item (Đúng) Phải xác định đúng các lực tác dụng và chọn chiều dương trước khi lập phương trình.
\item (Sai) Có thể bỏ qua điều kiện áp dụng và chiều của lực mà kết quả vẫn luôn đúng.
\end{itemchoice}
}
\end{ex}

% ===== Câu 148 =====
% ID: AIQ_L10C3_FULL_0958
% Mức: VD
\begin{ex}
% ID: AIQ_L10C3_FULL_0958
% Mức: VD
Muốn thùng $2\,\text{kg}$ trên sàn ngang chuyển động đều, hệ số ma sát trượt 0,5. Lấy g= $10\,\text{m}$ /s^2. Lực kéo ngang tối thiểu khi đã trượt là bao nhiêu?
\choiceTF
{\True Giá trị cần tìm là $10\,\text{N}$.}
{Giá trị cần tìm là $11\,\text{N}$.}
{\True Phải xác định đúng các lực tác dụng và chọn chiều dương trước khi lập phương trình.}
{Có thể bỏ qua điều kiện áp dụng và chiều của lực mà kết quả vẫn luôn đúng.}
\loigiai{
\begin{itemchoice}
\item (Đúng) Giá trị cần tìm là $10\,\text{N}$. (Vì): Chuyển động đều nên lực kéo cân bằng ma sát: F=μmg= $10\,\text{N}$. b) Sai vì sai giá trị. c) Đúng. d) Sai vì phải kiểm tra điều kiện áp dụng và chiều lực
\item (Sai) Giá trị cần tìm là $11\,\text{N}$.
\item (Đúng) Phải xác định đúng các lực tác dụng và chọn chiều dương trước khi lập phương trình.
\item (Sai) Có thể bỏ qua điều kiện áp dụng và chiều của lực mà kết quả vẫn luôn đúng.
\end{itemchoice}
}
\end{ex}

% ===== Câu 149 =====
% ID: AIQ_L10C3_FULL_0959
% Mức: VD
\begin{ex}
% ID: AIQ_L10C3_FULL_0959
% Mức: VD
Muốn thùng $3\,\text{kg}$ trên sàn ngang chuyển động đều, hệ số ma sát trượt 0,2. Lấy g= $10\,\text{m}$ /s^2. Lực kéo ngang tối thiểu khi đã trượt là bao nhiêu?
\choiceTF
{\True Giá trị cần tìm là $6\,\text{N}$.}
{Giá trị cần tìm là $7\,\text{N}$.}
{\True Phải xác định đúng các lực tác dụng và chọn chiều dương trước khi lập phương trình.}
{Có thể bỏ qua điều kiện áp dụng và chiều của lực mà kết quả vẫn luôn đúng.}
\loigiai{
\begin{itemchoice}
\item (Đúng) Giá trị cần tìm là $6\,\text{N}$. (Vì): Chuyển động đều nên lực kéo cân bằng ma sát: F=μmg= $6\,\text{N}$. b) Sai vì sai giá trị. c) Đúng. d) Sai vì phải kiểm tra điều kiện áp dụng và chiều lực
\item (Sai) Giá trị cần tìm là $7\,\text{N}$.
\item (Đúng) Phải xác định đúng các lực tác dụng và chọn chiều dương trước khi lập phương trình.
\item (Sai) Có thể bỏ qua điều kiện áp dụng và chiều của lực mà kết quả vẫn luôn đúng.
\end{itemchoice}
}
\end{ex}

% ===== Câu 150 =====
% ID: AIQ_L10C3_FULL_0960
% Mức: VD
\begin{ex}
% ID: AIQ_L10C3_FULL_0960
% Mức: VD
Muốn thùng $4\,\text{kg}$ trên sàn ngang chuyển động đều, hệ số ma sát trượt 0,3. Lấy g= $10\,\text{m}$ /s^2. Lực kéo ngang tối thiểu khi đã trượt là bao nhiêu?
\choiceTF
{\True Giá trị cần tìm là $12\,\text{N}$.}
{Giá trị cần tìm là $13\,\text{N}$.}
{\True Phải xác định đúng các lực tác dụng và chọn chiều dương trước khi lập phương trình.}
{Có thể bỏ qua điều kiện áp dụng và chiều của lực mà kết quả vẫn luôn đúng.}
\loigiai{
\begin{itemchoice}
\item (Đúng) Giá trị cần tìm là $12\,\text{N}$. (Vì): Chuyển động đều nên lực kéo cân bằng ma sát: F=μmg= $12\,\text{N}$. b) Sai vì sai giá trị. c) Đúng. d) Sai vì phải kiểm tra điều kiện áp dụng và chiều lực
\item (Sai) Giá trị cần tìm là $13\,\text{N}$.
\item (Đúng) Phải xác định đúng các lực tác dụng và chọn chiều dương trước khi lập phương trình.
\item (Sai) Có thể bỏ qua điều kiện áp dụng và chiều của lực mà kết quả vẫn luôn đúng.
\end{itemchoice}
}
\end{ex}

% ===== Câu 151 =====
% ID: AIQ_L10C3_FULL_0961
% Mức: TH
\begin{ex}
% ID: AIQ_L10C3_FULL_0961
% Mức: TH
Muốn thùng $3\,\text{kg}$ trên sàn ngang chuyển động đều, hệ số ma sát trượt 0,4. Lấy g= $10\,\text{m}$ /s^2. Lực kéo ngang tối thiểu khi đã trượt là bao nhiêu?
\shortans{12}
\loigiai{
Chuyển động đều nên lực kéo cân bằng ma sát: F=μmg= $12\,\text{N}$.
}
\end{ex}

% ===== Câu 152 =====
% ID: AIQ_L10C3_FULL_0962
% Mức: TH
\begin{ex}
% ID: AIQ_L10C3_FULL_0962
% Mức: TH
Muốn thùng $4\,\text{kg}$ trên sàn ngang chuyển động đều, hệ số ma sát trượt 0,5. Lấy g= $10\,\text{m}$ /s^2. Lực kéo ngang tối thiểu khi đã trượt là bao nhiêu?
\shortans{20}
\loigiai{
Chuyển động đều nên lực kéo cân bằng ma sát: F=μmg= $20\,\text{N}$.
}
\end{ex}

% ===== Câu 153 =====
% ID: AIQ_L10C3_FULL_0963
% Mức: VD
\begin{ex}
% ID: AIQ_L10C3_FULL_0963
% Mức: VD
Muốn thùng $5\,\text{kg}$ trên sàn ngang chuyển động đều, hệ số ma sát trượt 0,2. Lấy g= $10\,\text{m}$ /s^2. Lực kéo ngang tối thiểu khi đã trượt là bao nhiêu?
\shortans{10}
\loigiai{
Chuyển động đều nên lực kéo cân bằng ma sát: F=μmg= $10\,\text{N}$.
}
\end{ex}

% ===== Câu 154 =====
% ID: AIQ_L10C3_FULL_0964
% Mức: VD
\begin{ex}
% ID: AIQ_L10C3_FULL_0964
% Mức: VD
Muốn thùng $6\,\text{kg}$ trên sàn ngang chuyển động đều, hệ số ma sát trượt 0,3. Lấy g= $10\,\text{m}$ /s^2. Lực kéo ngang tối thiểu khi đã trượt là bao nhiêu?
\shortans{18}
\loigiai{
Chuyển động đều nên lực kéo cân bằng ma sát: F=μmg= $18\,\text{N}$.
}
\end{ex}

% ===== Câu 155 =====
% ID: AIQ_L10C3_FULL_0965
% Mức: VD
\begin{ex}
% ID: AIQ_L10C3_FULL_0965
% Mức: VD
Muốn thùng $2\,\text{kg}$ trên sàn ngang chuyển động đều, hệ số ma sát trượt 0,4. Lấy g= $10\,\text{m}$ /s^2. Lực kéo ngang tối thiểu khi đã trượt là bao nhiêu?
\shortans{8}
\loigiai{
Chuyển động đều nên lực kéo cân bằng ma sát: F=μmg= $8\,\text{N}$.
}
\end{ex}

% ===== Câu 156 =====
% ID: AIQ_L10C3_FULL_0966
% Mức: VDC
\begin{ex}
% ID: AIQ_L10C3_FULL_0966
% Mức: VDC
Muốn thùng $3\,\text{kg}$ trên sàn ngang chuyển động đều, hệ số ma sát trượt 0,5. Lấy g= $10\,\text{m}$ /s^2. Lực kéo ngang tối thiểu khi đã trượt là bao nhiêu?
\shortans{15}
\loigiai{
Chuyển động đều nên lực kéo cân bằng ma sát: F=μmg= $15\,\text{N}$.
}
\end{ex}

% ===== Câu 157 =====
% ID: AIQ_L10C3_FULL_0967
% Mức: VD
\begin{bt}
% ID: AIQ_L10C3_FULL_0967
% Mức: VD
Muốn thùng $6\,\text{kg}$ trên sàn ngang chuyển động đều, hệ số ma sát trượt 0,5. Lấy g= $10\,\text{m}$ /s^2. Lực kéo ngang tối thiểu khi đã trượt là bao nhiêu?
\loigiai{
Phân tích lực, chọn hệ quy chiếu và áp dụng định luật hoặc quy tắc phù hợp. Chuyển động đều nên lực kéo cân bằng ma sát: F=μmg= $30\,\text{N}$.
}
\end{bt}

% ===== Câu 158 =====
% ID: AIQ_L10C3_FULL_0968
% Mức: VD
\begin{bt}
% ID: AIQ_L10C3_FULL_0968
% Mức: VD
Muốn thùng $2\,\text{kg}$ trên sàn ngang chuyển động đều, hệ số ma sát trượt 0,2. Lấy g= $10\,\text{m}$ /s^2. Lực kéo ngang tối thiểu khi đã trượt là bao nhiêu?
\loigiai{
Phân tích lực, chọn hệ quy chiếu và áp dụng định luật hoặc quy tắc phù hợp. Chuyển động đều nên lực kéo cân bằng ma sát: F=μmg= $4\,\text{N}$.
}
\end{bt}

% ===== Câu 159 =====
% ID: AIQ_L10C3_FULL_0969
% Mức: VDC
\begin{bt}
% ID: AIQ_L10C3_FULL_0969
% Mức: VDC
Muốn thùng $3\,\text{kg}$ trên sàn ngang chuyển động đều, hệ số ma sát trượt 0,3. Lấy g= $10\,\text{m}$ /s^2. Lực kéo ngang tối thiểu khi đã trượt là bao nhiêu?
\loigiai{
Phân tích lực, chọn hệ quy chiếu và áp dụng định luật hoặc quy tắc phù hợp. Chuyển động đều nên lực kéo cân bằng ma sát: F=μmg= $9\,\text{N}$.
}
\end{bt}

% ===== Câu 160 =====
% ID: AIQ_L10C3_FULL_0970
% Mức: VDC
\begin{bt}
% ID: AIQ_L10C3_FULL_0970
% Mức: VDC
Muốn thùng $4\,\text{kg}$ trên sàn ngang chuyển động đều, hệ số ma sát trượt 0,4. Lấy g= $10\,\text{m}$ /s^2. Lực kéo ngang tối thiểu khi đã trượt là bao nhiêu?
\loigiai{
Phân tích lực, chọn hệ quy chiếu và áp dụng định luật hoặc quy tắc phù hợp. Chuyển động đều nên lực kéo cân bằng ma sát: F=μmg= $16\,\text{N}$.
}
\end{bt}

% ===== Câu 161 =====
% ID: AIQ_L10C3_FULL_0971
% Mức: VDC
\begin{bt}
% ID: AIQ_L10C3_FULL_0971
% Mức: VDC
Muốn thùng $5\,\text{kg}$ trên sàn ngang chuyển động đều, hệ số ma sát trượt 0,5. Lấy g= $10\,\text{m}$ /s^2. Lực kéo ngang tối thiểu khi đã trượt là bao nhiêu?
\loigiai{
Phân tích lực, chọn hệ quy chiếu và áp dụng định luật hoặc quy tắc phù hợp. Chuyển động đều nên lực kéo cân bằng ma sát: F=μmg= $25\,\text{N}$.
}
\end{bt}

% ===== Câu 162 =====
% ID: AIQ_L10C3_FULL_0972
% Mức: VDC
\begin{bt}
% ID: AIQ_L10C3_FULL_0972
% Mức: VDC
Muốn thùng $6\,\text{kg}$ trên sàn ngang chuyển động đều, hệ số ma sát trượt 0,2. Lấy g= $10\,\text{m}$ /s^2. Lực kéo ngang tối thiểu khi đã trượt là bao nhiêu?
\loigiai{
Phân tích lực, chọn hệ quy chiếu và áp dụng định luật hoặc quy tắc phù hợp. Chuyển động đều nên lực kéo cân bằng ma sát: F=μmg= $12\,\text{N}$.
}
\end{bt}
